
\Data\ID{1003076701}
\Updated{2020-07-15 03:07:12}
\Level{A}
\SubArea{Sine, cosine, tangent and cotangent}
\LANGen{
\Question{
The value of the expression \( 3\cos\frac{\pi}2 - 3\sin\pi + 2\left(\cos\frac{\pi}3 - \sin\frac{4\pi}3 \right) \) is:}
{\( 1+\sqrt3 \)}{\( 0  \)}{\( \frac12 \)}{\( 2 \)}{}{}}
\LANGpl{
\Question{
Wartość wyrażenia  \( 3\cos\frac{\pi}2 - 3\sin\pi + 2\left(\cos\frac{\pi}3 - \sin\frac{4\pi}3 \right) \) jest równa:}
{\( 1+\sqrt3 \)}{\( 0  \)}{\( \frac12 \)}{\( 2 \)}{}{}}
\LANGcs{
\Question{
Hodnota výrazu \( 3\cos\frac{\pi}2 - 3\sin\pi + 2\left(\cos\frac{\pi}3 - \sin\frac{4\pi}3 \right) \) je:}
{\( 1+\sqrt3 \)}{\( 0  \)}{\( \frac12 \)}{\( 2 \)}{}{}}
\LANGsk{
\Question{
Hodnota výrazu  \( 3\cos\frac{\pi}2 - 3\sin\pi + 2\left(\cos\frac{\pi}3 - \sin\frac{4\pi}3 \right) \) je:}
{\( 1+\sqrt3 \)}{\( 0  \)}{\( \frac12 \)}{\( 2 \)}{}{}}
\LANGes{
\Question{
El valor de la expresión \( 3\cos\frac{\pi}2 - 3\sin\pi + 2\left(\cos\frac{\pi}3 - \sin\frac{4\pi}3 \right) \) es:}
{\( 1+\sqrt3 \)}{\( 0  \)}{\( \frac12 \)}{\( 2 \)}{}{}}
\End

\Data\ID{9000032001}
\Updated{2020-07-15 03:07:34}
\Level{A}
\SubArea{Sine, cosine, tangent and cotangent}
\IMG{000320_01-figure0.pdf}
\LANGen{
\Question{
Evaluate
\(\mathop{\mathrm{tg}}\nolimits \left ( \frac{\pi }{2}\right )\).}
{undefined}{\(\frac{\sqrt{2}}
 {2} \)}{\(- 1\)}{\(\frac{\sqrt{3}}
 {3} \)}{\(-\frac{\sqrt{3}} 
 {3} \)}{\(-\sqrt{3}\)}}
\LANGpl{
\Question{
Oblicz
\(\mathop{\mathrm{tg}}\nolimits \left ( \frac{\pi }{2}\right )\).}
{nieokreślony}{\(\frac{\sqrt{2}}
 {2} \)}{\(- 1\)}{\(\frac{\sqrt{3}}
 {3} \)}{\(-\frac{\sqrt{3}} 
 {3} \)}{\(-\sqrt{3}\)}}
\LANGcs{
\Question{
\(\mathop{\mathrm{tg}}\nolimits \left ( \frac{\pi }{2}\right ) =?\)}
{není definován}{\(-\sqrt{3}\)}{\(- 1\)}{\(\frac{\sqrt{3}}3\)}{\(-\frac{\sqrt{3}}3\)}{\(\frac{\sqrt{2}}
 {2} \)}}
\LANGsk{
\Question{
\(\mathop{\mathrm{tg}}\nolimits \left ( \frac{\pi }{2}\right ) =?\)}
{nie je definované}{\(-\sqrt{3}\)}{\(- 1\)}{\(\frac{\sqrt{3}}3\)}{\(-\frac{\sqrt{3}}3\)}{\(\frac{\sqrt{2}}
 {2} \)}}
\LANGes{
\Question{
\(\mathop{\mathrm{tg}}\nolimits \left ( \frac{\pi }{2}\right ) = ?\)}
{no definido}{\(\frac{\sqrt{2}}
 {2} \)}{\(- 1\)}{\(\frac{\sqrt{3}}
 {3} \)}{\(-\frac{\sqrt{3}} 
 {3} \)}{\(-\sqrt{3}\)}}
\End

\Data\ID{9000032002}
\Updated{2020-07-15 03:07:34}
\Level{A}
\SubArea{Sine, cosine, tangent and cotangent}
\IMG{000320_02-figure0.pdf}
\LANGen{
\Question{
Evaluate
\(\mathop{\mathrm{tg}}\nolimits \left (0\right )\).}
{\(0\)}{\(\sqrt{3}\)}{\(\frac{\sqrt{3}}
 {3} \)}{\(\frac{\sqrt{2}}
 {2} \)}{\(1\)}{\(-\frac{\sqrt{3}} 
 {3} \)}}
\LANGpl{
\Question{
Oblicz
\(\mathop{\mathrm{tg}}\nolimits \left (0\right )\).}
{\(0\)}{\(\sqrt{3}\)}{\(\frac{\sqrt{3}}
 {3} \)}{\(\frac{\sqrt{2}}
 {2} \)}{\(1\)}{\(-\frac{\sqrt{3}} 
 {3} \)}}
\LANGcs{
\Question{
\(\mathop{\mathrm{tg}}\nolimits \left (0\right ) =?\)}
{\(0\)}{\(\frac{\sqrt{2}}
 {2} \)}{\(1\)}{\(-\sqrt{3}\)}{není definován}{\(\frac{\sqrt{3}}
 {3} \)}}
\LANGsk{
\Question{
\(\mathop{\mathrm{tg}}\nolimits \left (0\right ) =?\)}
{\(0\)}{\(\frac{\sqrt{2}}
 {2} \)}{\(1\)}{\(-\sqrt{3}\)}{nie je definované}{\(\frac{\sqrt{3}}
 {3} \)}}
\LANGes{
\Question{
\(\mathop{\mathrm{tg}}\nolimits \left (0\right ) = ?\)}
{\(0\)}{\(\sqrt{3}\)}{\(\frac{\sqrt{3}}
 {3} \)}{\(\frac{\sqrt{2}}
 {2} \)}{\(1\)}{\(-\frac{\sqrt{3}} 
 {3} \)}}
\End

\Data\ID{9000032003}
\Updated{2020-07-15 03:07:34}
\Level{A}
\SubArea{Sine, cosine, tangent and cotangent}
\IMG{000320_03-figure0.pdf}
\LANGen{
\Question{
Evaluate
\(\mathop{\mathrm{tg}}\nolimits \left (\frac{5\pi }
{2}\right )\).}
{undefined}{\(1\)}{\(- 1\)}{\(\sqrt{3}\)}{\(\frac{\sqrt{2}}
 {2} \)}{\(-\frac{\sqrt{3}} 
 {3} \)}}
\LANGpl{
\Question{
Oblicz
\(\mathop{\mathrm{tg}}\nolimits \left (\frac{5\pi }
{2}\right )\).}
{nieokreślony}{\(1\)}{\(- 1\)}{\(\sqrt{3}\)}{\(\frac{\sqrt{2}}
 {2} \)}{\(-\frac{\sqrt{3}} 
 {3} \)}}
\LANGcs{
\Question{
\(\mathop{\mathrm{tg}}\nolimits \left (\frac{5\pi }
{2}\right ) =?\)}
{není definován}{\(-\sqrt{3}\)}{\(-\frac{\sqrt{2}} 
 {2} \)}{\(0\)}{\(\frac{\sqrt{3}}
 {3} \)}{\(- 1\)}}
\LANGsk{
\Question{
\(\mathop{\mathrm{tg}}\nolimits \left (\frac{5\pi }
{2}\right ) =?\)}
{nie je definované}{\(-\sqrt{3}\)}{\(-\frac{\sqrt{2}} 
 {2} \)}{\(0\)}{\(\frac{\sqrt{3}}
 {3} \)}{\(- 1\)}}
\LANGes{
\Question{
\(\mathop{\mathrm{tg}}\nolimits \left (\frac{5\pi }
{2}\right ) = ?\)}
{no definido}{\(1\)}{\(- 1\)}{\(\sqrt{3}\)}{\(\frac{\sqrt{2}}
 {2} \)}{\(-\frac{\sqrt{3}} 
 {3} \)}}
\End

\Data\ID{9000032004}
\Updated{2020-07-15 03:07:34}
\Level{A}
\SubArea{Sine, cosine, tangent and cotangent}
\IMG{000320_04-figure0.pdf}
\LANGen{
\Question{
Evaluate
\(\mathop{\mathrm{tg}}\nolimits \left (-\frac{3\pi }
{2}\right )\).}
{undefined}{\(0\)}{\(1\)}{\(-\frac{\sqrt{3}} 
 {3} \)}{\(- 1\)}{\(\frac{\sqrt{3}}
 {3} \)}}
\LANGpl{
\Question{
Oblicz
\(\mathop{\mathrm{tg}}\nolimits \left (-\frac{3\pi }
{2}\right )\).}
{nieokreślony}{\(0\)}{\(1\)}{\(-\frac{\sqrt{3}} 
 {3} \)}{\(- 1\)}{\(\frac{\sqrt{3}}
 {3} \)}}
\LANGcs{
\Question{
\(\mathop{\mathrm{tg}}\nolimits \left (-\frac{3\pi }
{2}\right ) =?\)}
{není definován}{\(- 1\)}{\(1\)}{\(\frac{\sqrt{2}}
 {2} \)}{\(-\sqrt{3}\)}{\(\sqrt{3}\)}}
\LANGsk{
\Question{
\(\mathop{\mathrm{tg}}\nolimits \left (-\frac{3\pi }
{2}\right ) =?\)}
{nie je definované}{\(- 1\)}{\(1\)}{\(\frac{\sqrt{2}}
 {2} \)}{\(-\sqrt{3}\)}{\(\sqrt{3}\)}}
\LANGes{
\Question{
\(\mathop{\mathrm{tg}}\nolimits \left (-\frac{3\pi }
{2}\right ) = ?\)}
{no definido}{\(0\)}{\(1\)}{\(-\frac{\sqrt{3}} 
 {3} \)}{\(- 1\)}{\(\frac{\sqrt{3}}
 {3} \)}}
\End

\Data\ID{9000032005}
\Updated{2020-07-15 03:07:34}
\Level{A}
\SubArea{Sine, cosine, tangent and cotangent}
\IMG{000320_05-figure0.pdf}
\LANGen{
\Question{
Evaluate
\(\mathop{\mathrm{cotg}}\nolimits \left ( \frac{\pi }{2}\right )\).}
{\(0\)}{\(-\sqrt{3}\)}{\(-\frac{\sqrt{2}} 
 {2} \)}{\(\sqrt{3}\)}{\(\frac{\sqrt{2}}
 {2} \)}{\(1\)}}
\LANGpl{
\Question{
Oblicz
\(\mathop{\mathrm{cotg}}\nolimits \left ( \frac{\pi }{2}\right )\).}
{\(0\)}{\(-\sqrt{3}\)}{\(-\frac{\sqrt{2}} 
 {2} \)}{\(\sqrt{3}\)}{\(\frac{\sqrt{2}}
 {2} \)}{\(1\)}}
\LANGcs{
\Question{
\(\mathop{\mathrm{cotg}}\nolimits \left ( \frac{\pi }{2}\right ) =?\)}
{\(0\)}{\(\frac{\sqrt{3}}
 {3} \)}{\(\frac{\sqrt{2}}
 {2} \)}{\(1\)}{\(-\frac{\sqrt{2}} 
 {2} \)}{není definován}}
\LANGsk{
\Question{
\(\mathop{\mathrm{cotg}}\nolimits \left ( \frac{\pi }{2}\right ) =?\)}
{\(0\)}{\(\frac{\sqrt{3}}
 {3} \)}{\(\frac{\sqrt{2}}
 {2} \)}{\(1\)}{\(-\frac{\sqrt{2}} 
 {2} \)}{nie je definované}}
\LANGes{
\Question{
\(\mathop{\mathrm{cotg}}\nolimits \left ( \frac{\pi }{2}\right ) = ?\)}
{\(0\)}{\(-\sqrt{3}\)}{\(-\frac{\sqrt{2}} 
 {2} \)}{\(\sqrt{3}\)}{\(\frac{\sqrt{2}}
 {2} \)}{\(1\)}}
\End

\Data\ID{9000032006}
\Updated{2020-07-15 03:07:34}
\Level{A}
\SubArea{Sine, cosine, tangent and cotangent}
\IMG{000320_06-figure0.pdf}
\LANGen{
\Question{
Evaluate
\(\mathop{\mathrm{cotg}}\nolimits \left (0\right )\).}
{undefined}{\(-\sqrt{3}\)}{\(-\frac{\sqrt{2}} 
 {2} \)}{\(\frac{\sqrt{3}}
 {3} \)}{\(-\frac{\sqrt{3}} 
 {3} \)}{\(\frac{\sqrt{2}}
 {2} \)}}
\LANGpl{
\Question{
Oblicz
\(\mathop{\mathrm{cotg}}\nolimits \left (0\right )\).}
{nieokreślony}{\(-\sqrt{3}\)}{\(-\frac{\sqrt{2}} 
 {2} \)}{\(\frac{\sqrt{3}}
 {3} \)}{\(-\frac{\sqrt{3}} 
 {3} \)}{\(\frac{\sqrt{2}}
 {2} \)}}
\LANGcs{
\Question{
\(\mathop{\mathrm{cotg}}\nolimits \left (0\right ) =?\)}
{není definován}{\(\frac{\sqrt{2}}
 {2} \)}{\(\sqrt{3}\)}{\(0\)}{\(-\sqrt{3}\)}{\(-\frac{\sqrt{2}} 
 {2} \)}}
\LANGsk{
\Question{
\(\mathop{\mathrm{cotg}}\nolimits \left (0\right ) =?\)}
{nie je definované}{\(\frac{\sqrt{2}}
 {2} \)}{\(\sqrt{3}\)}{\(0\)}{\(-\sqrt{3}\)}{\(-\frac{\sqrt{2}} 
 {2} \)}}
\LANGes{
\Question{
\(\mathop{\mathrm{cotg}}\nolimits \left (0\right ) = ?\)}
{no definido}{\(-\sqrt{3}\)}{\(-\frac{\sqrt{2}} 
 {2} \)}{\(\frac{\sqrt{3}}
 {3} \)}{\(-\frac{\sqrt{3}} 
 {3} \)}{\(\frac{\sqrt{2}}
 {2} \)}}
\End

\Data\ID{9000032007}
\Updated{2020-07-15 03:07:34}
\Level{A}
\SubArea{Sine, cosine, tangent and cotangent}
\IMG{000320_07-figure0.pdf}
\LANGen{
\Question{
Evaluate
\(\mathop{\mathrm{cotg}}\nolimits \left (\frac{5\pi }
{2}\right )\).}
{\(0\)}{\(\sqrt{3}\)}{\(-\frac{\sqrt{2}} 
 {2} \)}{\(-\frac{\sqrt{3}} 
 {3} \)}{\(\frac{\sqrt{3}}
 {3} \)}{undefined}}
\LANGpl{
\Question{
Oblicz
\(\mathop{\mathrm{cotg}}\nolimits \left (\frac{5\pi }
{2}\right )\).}
{\(0\)}{\(\sqrt{3}\)}{\(-\frac{\sqrt{2}} 
 {2} \)}{\(-\frac{\sqrt{3}} 
 {3} \)}{\(\frac{\sqrt{3}}
 {3} \)}{nieokreślony}}
\LANGcs{
\Question{
\(\mathop{\mathrm{cotg}}\nolimits \left (\frac{5\pi }
{2}\right ) =?\)}
{\(0\)}{\(\sqrt{3}\)}{\(1\)}{není definován}{\(-\frac{\sqrt{2}} 
 {2} \)}{\(\frac{\sqrt{2}}
 {2} \)}}
\LANGsk{
\Question{
\(\mathop{\mathrm{cotg}}\nolimits \left (\frac{5\pi }
{2}\right ) =?\)}
{\(0\)}{\(\sqrt{3}\)}{\(1\)}{nie je definované}{\(-\frac{\sqrt{2}} 
 {2} \)}{\(\frac{\sqrt{2}}
 {2} \)}}
\LANGes{
\Question{
\(\mathop{\mathrm{cotg}}\nolimits \left (\frac{5\pi }
{2}\right) = ?\)}
{\(0\)}{\(\sqrt{3}\)}{\(-\frac{\sqrt{2}} 
 {2} \)}{\(-\frac{\sqrt{3}} 
 {3} \)}{\(\frac{\sqrt{3}}
 {3} \)}{no definido}}
\End

\Data\ID{9000032008}
\Updated{2020-07-15 03:07:34}
\Level{A}
\SubArea{Sine, cosine, tangent and cotangent}
\IMG{000320_08-figure0.pdf}
\LANGen{
\Question{
Evaluate
\(\mathop{\mathrm{cotg}}\nolimits \left (-\frac{3\pi }
{2}\right )\).}
{\(0\)}{\(\frac{\sqrt{2}}
 {2} \)}{\(\sqrt{3}\)}{undefined}{\(1\)}{\(-\frac{\sqrt{2}} 
 {2} \)}}
\LANGpl{
\Question{
Oblicz
\(\mathop{\mathrm{cotg}}\nolimits \left (-\frac{3\pi }
{2}\right )\).}
{\(0\)}{\(\frac{\sqrt{2}}
 {2} \)}{\(\sqrt{3}\)}{nieokreślony}{\(1\)}{\(-\frac{\sqrt{2}} 
 {2} \)}}
\LANGcs{
\Question{
\(\mathop{\mathrm{cotg}}\nolimits \left (-\frac{3\pi }
{2}\right ) =?\)}
{\(0\)}{\(-\sqrt{3}\)}{\(\frac{\sqrt{3}}
 {3} \)}{\(\sqrt{3}\)}{\(-\frac{\sqrt{3}} 
 {3} \)}{není definován}}
\LANGsk{
\Question{
\(\mathop{\mathrm{cotg}}\nolimits \left (-\frac{3\pi }
{2}\right ) =?\)}
{\(0\)}{\(-\sqrt{3}\)}{\(\frac{\sqrt{3}}
 {3} \)}{\(\sqrt{3}\)}{\(-\frac{\sqrt{3}} 
 {3} \)}{nie je definované}}
\LANGes{
\Question{
\(\mathop{\mathrm{cotg}}\nolimits \left (-\frac{3\pi }
{2}\right ) = ?\)}
{\(0\)}{\(\frac{\sqrt{2}}
 {2} \)}{\(\sqrt{3}\)}{no definido}{\(1\)}{\(-\frac{\sqrt{2}} 
 {2} \)}}
\End

\Data\ID{9000032009}
\Updated{2020-07-15 03:07:34}
\Level{A}
\SubArea{Sine, cosine, tangent and cotangent}
\IMG{000320_09-figure0.pdf}
\LANGen{
\Question{
Evaluate
\(\mathop{\mathrm{tg}}\nolimits \left ( \frac{\pi }{4}\right )\).}
{\(1\)}{\(\frac{\sqrt{3}}
 {3} \)}{\(0\)}{\(\frac{\sqrt{2}}
 {2} \)}{\(-\frac{\sqrt{2}} 
 {2} \)}{\(- 1\)}}
\LANGpl{
\Question{
Oblicz
\(\mathop{\mathrm{tg}}\nolimits \left ( \frac{\pi }{4}\right )\).}
{\(1\)}{\(\frac{\sqrt{3}}
 {3} \)}{\(0\)}{\(\frac{\sqrt{2}}
 {2} \)}{\(-\frac{\sqrt{2}} 
 {2} \)}{\(- 1\)}}
\LANGcs{
\Question{
\(\mathop{\mathrm{tg}}\nolimits \left ( \frac{\pi }{4}\right ) =?\)}
{\(1\)}{\(\frac{\sqrt{3}}
 {3} \)}{\(\sqrt{3}\)}{není definován}{\(-\frac{\sqrt{2}} 
 {2} \)}{\(-\frac{\sqrt{3}} 
 {3} \)}}
\LANGsk{
\Question{
\(\mathop{\mathrm{tg}}\nolimits \left ( \frac{\pi }{4}\right ) =?\)}
{\(1\)}{\(\frac{\sqrt{3}}
 {3} \)}{\(\sqrt{3}\)}{nie je definované}{\(-\frac{\sqrt{2}} 
 {2} \)}{\(-\frac{\sqrt{3}} 
 {3} \)}}
\LANGes{
\Question{
\(\mathop{\mathrm{tg}}\nolimits \left ( \frac{\pi }{4}\right) = ?\)}
{\(1\)}{\(\frac{\sqrt{3}}
 {3} \)}{\(0\)}{\(\frac{\sqrt{2}}
 {2} \)}{\(-\frac{\sqrt{2}} 
 {2} \)}{\(- 1\)}}
\End

\Data\ID{9000032010}
\Updated{2020-07-15 03:07:34}
\Level{A}
\SubArea{Sine, cosine, tangent and cotangent}
\IMG{000320_10-figure0.pdf}
\LANGen{
\Question{
Evaluate
\(\mathop{\mathrm{cotg}}\nolimits \left ( \frac{\pi }{4}\right )\).}
{\(1\)}{\(\frac{\sqrt{2}}
 {2} \)}{\(0\)}{\(-\frac{\sqrt{3}} 
 {3} \)}{\(\frac{\sqrt{3}}
 {3} \)}{undefined}}
\LANGpl{
\Question{
Oblicz
\(\mathop{\mathrm{cotg}}\nolimits \left ( \frac{\pi }{4}\right )\).}
{\(1\)}{\(\frac{\sqrt{2}}
 {2} \)}{\(0\)}{\(-\frac{\sqrt{3}} 
 {3} \)}{\(\frac{\sqrt{3}}
 {3} \)}{nieokreślony}}
\LANGcs{
\Question{
\(\mathop{\mathrm{cotg}}\nolimits \left ( \frac{\pi }{4}\right ) =?\)}
{\(1\)}{\(-\frac{\sqrt{3}} 
 {3} \)}{\(-\frac{\sqrt{2}} 
 {2} \)}{není definován}{\(0\)}{\(\frac{\sqrt{3}}
 {3} \)}}
\LANGsk{
\Question{
\(\mathop{\mathrm{cotg}}\nolimits \left ( \frac{\pi }{4}\right ) =?\)}
{\(1\)}{\(-\frac{\sqrt{3}} 
 {3} \)}{\(-\frac{\sqrt{2}} 
 {2} \)}{nie je definované}{\(0\)}{\(\frac{\sqrt{3}}
 {3} \)}}
\LANGes{
\Question{
\(\mathop{\mathrm{cotg}}\nolimits \left ( \frac{\pi }{4}\right ) = ?\)}
{\(1\)}{\(\frac{\sqrt{2}}
 {2} \)}{\(0\)}{\(-\frac{\sqrt{3}} 
 {3} \)}{\(\frac{\sqrt{3}}
 {3} \)}{no definido}}
\End

\Data\ID{9000032011}
\Updated{2020-07-15 03:07:34}
\Level{A}
\SubArea{Sine, cosine, tangent and cotangent}
\IMG{000320_11-figure0.pdf}
\LANGen{
\Question{
Evaluate
\(\mathop{\mathrm{tg}}\nolimits \left ( \frac{\pi }{3}\right )\).}
{\(\sqrt{3}\)}{\(0\)}{\(-\sqrt{3}\)}{\(-\frac{\sqrt{3}} 
 {3} \)}{\(-\frac{\sqrt{2}} 
 {2} \)}{\(\frac{\sqrt{3}}
 {3} \)}}
\LANGpl{
\Question{
Oblicz
\(\mathop{\mathrm{tg}}\nolimits \left ( \frac{\pi }{3}\right )\).}
{\(\sqrt{3}\)}{\(0\)}{\(-\sqrt{3}\)}{\(-\frac{\sqrt{3}} 
 {3} \)}{\(-\frac{\sqrt{2}} 
 {2} \)}{\(\frac{\sqrt{3}}
 {3} \)}}
\LANGcs{
\Question{
\(\mathop{\mathrm{tg}}\nolimits \left ( \frac{\pi }{3}\right ) =?\)}
{\(\sqrt{3}\)}{\(\frac{1}
{2}\)}{\(-\sqrt{3}\)}{\(\frac{\sqrt{3}}
 {3} \)}{\(0\)}{\(-\frac{1} 
{2}\)}}
\LANGsk{
\Question{
\(\mathop{\mathrm{tg}}\nolimits \left ( \frac{\pi }{3}\right ) =?\)}
{\(\sqrt{3}\)}{\(\frac{1}
{2}\)}{\(-\sqrt{3}\)}{\(\frac{\sqrt{3}}
 {3} \)}{\(0\)}{\(-\frac{1} 
{2}\)}}
\LANGes{
\Question{
\(\mathop{\mathrm{tg}}\nolimits \left ( \frac{\pi }{3}\right ) = ?\).}
{\(\sqrt{3}\)}{\(0\)}{\(-\sqrt{3}\)}{\(-\frac{\sqrt{3}} 
 {3} \)}{\(-\frac{\sqrt{2}} 
 {2} \)}{\(\frac{\sqrt{3}}
 {3} \)}}
\End

\Data\ID{9000032012}
\Updated{2020-07-15 04:07:14}
\Level{A}
\SubArea{Sine, cosine, tangent and cotangent}
\IMG{000320_12-figure0.pdf}
\LANGen{
\Question{
Evaluate
\(\mathop{\mathrm{cotg}}\nolimits \left ( \frac{\pi }{3}\right )\).}
{\(\frac{\sqrt{3}}
 {3} \)}{\(-\frac{1} 
{2}\)}{\(\frac{\sqrt{2}}
 {2} \)}{\(-\sqrt{3}\)}{\(-\frac{\sqrt{3}} 
 {3} \)}{\(-\frac{\sqrt{2}} 
 {2} \)}}
\LANGpl{
\Question{
Oblicz
\(\mathop{\mathrm{cotg}}\nolimits \left ( \frac{\pi }{3}\right )\).}
{\(\frac{\sqrt{3}}
 {3} \)}{\(-\frac{1} 
{2}\)}{\(\frac{\sqrt{2}}
 {2} \)}{\(-\sqrt{3}\)}{\(-\frac{\sqrt{3}} 
 {3} \)}{\(-\frac{\sqrt{2}} 
 {2} \)}}
\LANGcs{
\Question{
\(\mathop{\mathrm{cotg}}\nolimits \left ( \frac{\pi }{3}\right ) =?\)}
{\(\frac{\sqrt{3}}
 {3} \)}{\(0\)}{\(-\sqrt{3}\)}{\(-\frac{\sqrt{2}} 
 {2} \)}{\(\frac{1}
{2}\)}{\(-\frac{\sqrt{3}} 
 {3} \)}}
\LANGsk{
\Question{
\(\mathop{\mathrm{cotg}}\nolimits \left ( \frac{\pi }{3}\right ) =?\)}
{\(\frac{\sqrt{3}}
 {3} \)}{\(0\)}{\(-\sqrt{3}\)}{\(-\frac{\sqrt{2}} 
 {2} \)}{\(\frac{1}
{2}\)}{\(-\frac{\sqrt{3}} 
 {3} \)}}
\LANGes{
\Question{
\(\mathop{\mathrm{cotg}}\nolimits \left ( \frac{\pi }{3}\right ) = ?\)}
{\(\frac{\sqrt{3}}
 {3} \)}{\(-\frac{1} 
{2}\)}{\(\frac{\sqrt{2}}
 {2} \)}{\(-\sqrt{3}\)}{\(-\frac{\sqrt{3}} 
 {3} \)}{\(-\frac{\sqrt{2}} 
 {2} \)}}
\End

\Data\ID{9000032013}
\Updated{2020-07-15 03:07:34}
\Level{A}
\SubArea{Sine, cosine, tangent and cotangent}
\IMG{000320_13-figure0.pdf}
\LANGen{
\Question{
Evaluate
\(\mathop{\mathrm{tg}}\nolimits \left ( \frac{\pi }{6}\right )\).}
{\(\frac{\sqrt{3}}
 {3} \)}{\(0\)}{\(-\sqrt{3}\)}{\(-\frac{\sqrt{3}} 
 {3} \)}{\(\frac{\sqrt{2}}
 {2} \)}{\(-\frac{\sqrt{2}} 
 {2} \)}}
\LANGpl{
\Question{
Oblicz
\(\mathop{\mathrm{tg}}\nolimits \left ( \frac{\pi }{6}\right )\).}
{\(\frac{\sqrt{3}}
 {3} \)}{\(0\)}{\(-\sqrt{3}\)}{\(-\frac{\sqrt{3}} 
 {3} \)}{\(\frac{\sqrt{2}}
 {2} \)}{\(-\frac{\sqrt{2}} 
 {2} \)}}
\LANGcs{
\Question{
\(\mathop{\mathrm{tg}}\nolimits \left ( \frac{\pi }{6}\right ) =?\)}
{\(\frac{\sqrt{3}}
 {3} \)}{\(-\frac{1} 
{2}\)}{\(-\frac{\sqrt{3}} 
 {3} \)}{\(\frac{\sqrt{2}}
 {2} \)}{není definován}{\(\frac{1}
{2}\)}}
\LANGsk{
\Question{
\(\mathop{\mathrm{tg}}\nolimits \left ( \frac{\pi }{6}\right ) =?\)}
{\(\frac{\sqrt{3}}
 {3} \)}{\(-\frac{1} 
{2}\)}{\(-\frac{\sqrt{3}} 
 {3} \)}{\(\frac{\sqrt{2}}
 {2} \)}{nie je definované}{\(\frac{1}
{2}\)}}
\LANGes{
\Question{
\(\mathop{\mathrm{tg}}\nolimits \left ( \frac{\pi }{6}\right ) = ?\)}
{\(\frac{\sqrt{3}}
 {3} \)}{\(0\)}{\(-\sqrt{3}\)}{\(-\frac{\sqrt{3}} 
 {3} \)}{\(\frac{\sqrt{2}}
 {2} \)}{\(-\frac{\sqrt{2}} 
 {2} \)}}
\End

\Data\ID{9000032014}
\Updated{2020-07-15 03:07:34}
\Level{A}
\SubArea{Sine, cosine, tangent and cotangent}
\IMG{000320_14-figure0.pdf}
\LANGen{
\Question{
Evaluate
\(\mathop{\mathrm{cotg}}\nolimits \left ( \frac{\pi }{6}\right )\).}
{\(\sqrt{3}\)}{\(-\frac{1} 
{2}\)}{\(-\frac{\sqrt{3}} 
 {2} \)}{\(0\)}{undefined}{\(-\sqrt{3}\)}}
\LANGpl{
\Question{
Oblicz
\(\mathop{\mathrm{cotg}}\nolimits \left ( \frac{\pi }{6}\right )\).}
{\(\sqrt{3}\)}{\(-\frac{1} 
{2}\)}{\(-\frac{\sqrt{3}} 
 {2} \)}{\(0\)}{nieokreślony}{\(-\sqrt{3}\)}}
\LANGcs{
\Question{
\(\mathop{\mathrm{cotg}}\nolimits \left ( \frac{\pi }{6}\right ) =?\)}
{\(\sqrt{3}\)}{není definován}{\(-\frac{\sqrt{2}} 
 {2} \)}{\(\frac{1}
{2}\)}{\(0\)}{\(-\frac{\sqrt{3}} 
 {2} \)}}
\LANGsk{
\Question{
\(\mathop{\mathrm{cotg}}\nolimits \left ( \frac{\pi }{6}\right ) =?\)}
{\(\sqrt{3}\)}{nie je definované}{\(-\frac{\sqrt{2}} 
 {2} \)}{\(\frac{1}
{2}\)}{\(0\)}{\(-\frac{\sqrt{3}} 
 {2} \)}}
\LANGes{
\Question{
\(\mathop{\mathrm{cotg}}\nolimits \left ( \frac{\pi }{6}\right ) = ?\)}
{\(\sqrt{3}\)}{\(-\frac{1} 
{2}\)}{\(-\frac{\sqrt{3}} 
 {2} \)}{\(0\)}{no definido}{\(-\sqrt{3}\)}}
\End

\Data\ID{9000032101}
\Updated{2020-07-15 03:07:34}
\Level{A}
\SubArea{Sine, cosine, tangent and cotangent}
\IMG{000321_01-figure0.pdf}
\LANGen{
\Question{
Evaluate
\(\sin \left ( \frac{\pi }{2}\right )\).}
{\(1\)}{\(-\frac{\sqrt{2}} 
 {2} \)}{\(-\sqrt{3}\)}{\(\sqrt{3}\)}{\(\frac{\sqrt{2}}
 {2} \)}{\(\frac{\sqrt{3}}
 {3} \)}}
\LANGpl{
\Question{
Oblicz
\(\sin \left ( \frac{\pi }{2}\right )\).}
{\(1\)}{\(-\frac{\sqrt{2}} 
 {2} \)}{\(-\sqrt{3}\)}{\(\sqrt{3}\)}{\(\frac{\sqrt{2}}
 {2} \)}{\(\frac{\sqrt{3}}
 {3} \)}}
\LANGcs{
\Question{
\(\sin \left ( \frac{\pi }{2}\right ) =?\)}
{\(1\)}{\(-\sqrt{3}\)}{\(- 1\)}{\(0\)}{\(-\frac{\sqrt{2}} 
 {2} \)}{\(\frac{\sqrt{3}}
 {3} \)}}
\LANGsk{
\Question{
\(\sin \left ( \frac{\pi }{2}\right ) =?\)}
{\(1\)}{\(-\sqrt{3}\)}{\(- 1\)}{\(0\)}{\(-\frac{\sqrt{2}} 
 {2} \)}{\(\frac{\sqrt{3}}
 {3} \)}}
\LANGes{
\Question{
\(\sin \left ( \frac{\pi }{2}\right ) = ?\)}
{\(1\)}{\(-\frac{\sqrt{2}} 
 {2} \)}{\(-\sqrt{3}\)}{\(\sqrt{3}\)}{\(\frac{\sqrt{2}}
 {2} \)}{\(\frac{\sqrt{3}}
 {3} \)}}
\End

\Data\ID{9000032102}
\Updated{2020-07-15 03:07:34}
\Level{A}
\SubArea{Sine, cosine, tangent and cotangent}
\IMG{000321_02-figure0.pdf}
\LANGen{
\Question{
Evaluate
\(\sin \left (0\right )\).}
{\(0\)}{\(\frac{\sqrt{2}}
 {2} \)}{\(-\sqrt{3}\)}{\(- 1\)}{\(\sqrt{3}\)}{\(-\frac{\sqrt{3}} 
 {3} \)}}
\LANGpl{
\Question{
Oblicz
\(\sin \left (0\right )\).}
{\(0\)}{\(\frac{\sqrt{2}}
 {2} \)}{\(-\sqrt{3}\)}{\(- 1\)}{\(\sqrt{3}\)}{\(-\frac{\sqrt{3}} 
 {3} \)}}
\LANGcs{
\Question{
\(\sin \left (0\right ) =?\)}
{\(0\)}{\(\sqrt{3}\)}{\(-\frac{\sqrt{2}} 
 {2} \)}{\(-\frac{\sqrt{3}} 
 {3} \)}{\(1\)}{\(\frac{\sqrt{2}}
 {2} \)}}
\LANGsk{
\Question{
\(\sin \left (0\right ) =?\)}
{\(0\)}{\(\sqrt{3}\)}{\(-\frac{\sqrt{2}} 
 {2} \)}{\(-\frac{\sqrt{3}} 
 {3} \)}{\(1\)}{\(\frac{\sqrt{2}}
 {2} \)}}
\LANGes{
\Question{
\(\sin \left (0\right ) = ?\)}
{\(0\)}{\(\frac{\sqrt{2}}
 {2} \)}{\(-\sqrt{3}\)}{\(- 1\)}{\(\sqrt{3}\)}{\(-\frac{\sqrt{3}} 
 {3} \)}}
\End

\Data\ID{9000032103}
\Updated{2020-07-15 03:07:34}
\Level{A}
\SubArea{Sine, cosine, tangent and cotangent}
\IMG{000321_03-figure0.pdf}
\LANGen{
\Question{
Evaluate
\(\sin \left (\frac{5\pi }
{2}\right )\).}
{\(1\)}{\(-\frac{\sqrt{3}} 
 {3} \)}{\(-\frac{\sqrt{2}} 
 {2} \)}{\(\sqrt{3}\)}{\(\frac{\sqrt{3}}
 {3} \)}{\(\frac{\sqrt{2}}
 {2} \)}}
\LANGpl{
\Question{
Oblicz
\(\sin \left (\frac{5\pi }
{2}\right )\).}
{\(1\)}{\(-\frac{\sqrt{3}} 
 {3} \)}{\(-\frac{\sqrt{2}} 
 {2} \)}{\(\sqrt{3}\)}{\(\frac{\sqrt{3}}
 {3} \)}{\(\frac{\sqrt{2}}
 {2} \)}}
\LANGcs{
\Question{
\(\sin \left (\frac{5\pi }
{2}\right ) =?\)}
{\(1\)}{\(- 1\)}{\(-\sqrt{3}\)}{\(-\frac{\sqrt{2}} 
 {2} \)}{\(0\)}{\(\sqrt{3}\)}}
\LANGsk{
\Question{
\(\sin \left (\frac{5\pi }
{2}\right ) =?\)}
{\(1\)}{\(- 1\)}{\(-\sqrt{3}\)}{\(-\frac{\sqrt{2}} 
 {2} \)}{\(0\)}{\(\sqrt{3}\)}}
\LANGes{
\Question{
\(\sin \left (\frac{5\pi }
{2}\right ) = ?\)}
{\(1\)}{\(-\frac{\sqrt{3}} 
 {3} \)}{\(-\frac{\sqrt{2}} 
 {2} \)}{\(\sqrt{3}\)}{\(\frac{\sqrt{3}}
 {3} \)}{\(\frac{\sqrt{2}}
 {2} \)}}
\End

\Data\ID{9000032104}
\Updated{2020-07-15 03:07:34}
\Level{A}
\SubArea{Sine, cosine, tangent and cotangent}
\IMG{000321_04-figure0.pdf}
\LANGen{
\Question{
Evaluate
\(\sin \left (\frac{-3\pi }
 {2} \right )\).}
{\(1\)}{\(\frac{\sqrt{3}}
 {3} \)}{\(\sqrt{3}\)}{\(-\sqrt{3}\)}{\(-\frac{\sqrt{2}} 
 {2} \)}{\(-\frac{\sqrt{3}} 
 {3} \)}}
\LANGpl{
\Question{
Oblicz
\(\sin \left (\frac{-3\pi }
 {2} \right )\).}
{\(1\)}{\(\frac{\sqrt{3}}
 {3} \)}{\(\sqrt{3}\)}{\(-\sqrt{3}\)}{\(-\frac{\sqrt{2}} 
 {2} \)}{\(-\frac{\sqrt{3}} 
 {3} \)}}
\LANGcs{
\Question{
\(\sin \left (\frac{-3\pi }
 {2} \right ) =?\)}
{\(1\)}{\(\frac{\sqrt{3}}
 {3} \)}{\(\sqrt{3}\)}{\(0\)}{\(\frac{\sqrt{2}}
 {2} \)}{\(-\sqrt{3}\)}}
\LANGsk{
\Question{
\(\sin \left (\frac{-3\pi }
 {2} \right ) =?\)}
{\(1\)}{\(\frac{\sqrt{3}}
 {3} \)}{\(\sqrt{3}\)}{\(0\)}{\(\frac{\sqrt{2}}
 {2} \)}{\(-\sqrt{3}\)}}
\LANGes{
\Question{
\(\sin \left (\frac{-3\pi }
 {2} \right ) = ?\)}
{\(1\)}{\(\frac{\sqrt{3}}
 {3} \)}{\(\sqrt{3}\)}{\(-\sqrt{3}\)}{\(-\frac{\sqrt{2}} 
 {2} \)}{\(-\frac{\sqrt{3}} 
 {3} \)}}
\End

\Data\ID{9000032105}
\Updated{2020-07-15 03:07:34}
\Level{A}
\SubArea{Sine, cosine, tangent and cotangent}
\IMG{000321_05-figure0.pdf}
\LANGen{
\Question{
Evaluate
\(\cos \left ( \frac{\pi }{2}\right )\).}
{\(0\)}{\(- 1\)}{\(\frac{\sqrt{3}}
 {3} \)}{\(\sqrt{3}\)}{\(1\)}{\(-\frac{\sqrt{3}} 
 {3} \)}}
\LANGpl{
\Question{
Oblicz
\(\cos \left ( \frac{\pi }{2}\right )\).}
{\(0\)}{\(- 1\)}{\(\frac{\sqrt{3}}
 {3} \)}{\(\sqrt{3}\)}{\(1\)}{\(-\frac{\sqrt{3}} 
 {3} \)}}
\LANGcs{
\Question{
\(\cos \left ( \frac{\pi }{2}\right ) =?\)}
{\(0\)}{\(-\frac{\sqrt{2}} 
 {2} \)}{\(\frac{\sqrt{2}}
 {2} \)}{\(- 1\)}{\(1\)}{\(-\sqrt{3}\)}}
\LANGsk{
\Question{
\(\cos \left ( \frac{\pi }{2}\right ) =?\)}
{\(0\)}{\(-\frac{\sqrt{2}} 
 {2} \)}{\(\frac{\sqrt{2}}
 {2} \)}{\(- 1\)}{\(1\)}{\(-\sqrt{3}\)}}
\LANGes{
\Question{
\(\cos \left ( \frac{\pi }{2}\right) = ?\)}
{\(0\)}{\(- 1\)}{\(\frac{\sqrt{3}}
 {3} \)}{\(\sqrt{3}\)}{\(1\)}{\(-\frac{\sqrt{3}} 
 {3} \)}}
\End

\Data\ID{9000032106}
\Updated{2020-07-15 03:07:34}
\Level{A}
\SubArea{Sine, cosine, tangent and cotangent}
\IMG{000321_06-figure0.pdf}
\LANGen{
\Question{
Evaluate
\(\cos \left (0\right )\).}
{\(1\)}{\(\frac{\sqrt{3}}
 {3} \)}{\(\sqrt{3}\)}{\(- 1\)}{\(-\frac{\sqrt{3}} 
 {3} \)}{\(0\)}}
\LANGpl{
\Question{
Oblicz
\(\cos \left (0\right )\).}
{\(1\)}{\(\frac{\sqrt{3}}
 {3} \)}{\(\sqrt{3}\)}{\(- 1\)}{\(-\frac{\sqrt{3}} 
 {3} \)}{\(0\)}}
\LANGcs{
\Question{
\(\cos \left (0\right ) =?\)}
{\(1\)}{\(0\)}{\(- 1\)}{\(-\sqrt{3}\)}{\(\sqrt{3}\)}{\(-\frac{\sqrt{2}} 
 {2} \)}}
\LANGsk{
\Question{
\(\cos \left (0\right ) =?\)}
{\(1\)}{\(0\)}{\(- 1\)}{\(-\sqrt{3}\)}{\(\sqrt{3}\)}{\(-\frac{\sqrt{2}} 
 {2} \)}}
\LANGes{
\Question{
\(\cos \left (0\right ) = ?\)}
{\(1\)}{\(\frac{\sqrt{3}}
 {3} \)}{\(\sqrt{3}\)}{\(- 1\)}{\(-\frac{\sqrt{3}} 
 {3} \)}{\(0\)}}
\End

\Data\ID{9000032107}
\Updated{2020-07-15 03:07:34}
\Level{A}
\SubArea{Sine, cosine, tangent and cotangent}
\IMG{000321_07-figure0.pdf}
\LANGen{
\Question{
Evaluate
\(\cos \left (\pi \right )\).}
{\(- 1\)}{\(\sqrt{3}\)}{\(\frac{\sqrt{3}}
 {3} \)}{\(-\sqrt{3}\)}{\(-\frac{\sqrt{3}} 
 {3} \)}{\(0\)}}
\LANGpl{
\Question{
Oblicz
\(\cos \left (\pi \right )\).}
{\(- 1\)}{\(\sqrt{3}\)}{\(\frac{\sqrt{3}}
 {3} \)}{\(-\sqrt{3}\)}{\(-\frac{\sqrt{3}} 
 {3} \)}{\(0\)}}
\LANGcs{
\Question{
\(\cos \left (\pi \right ) =?\)}
{\(- 1\)}{\(\frac{\sqrt{2}}
 {2} \)}{\(0\)}{\(1\)}{\(-\frac{\sqrt{3}} 
 {3} \)}{\(-\frac{\sqrt{2}} 
 {2} \)}}
\LANGsk{
\Question{
\(\cos \left (\pi \right ) =?\)}
{\(- 1\)}{\(\frac{\sqrt{2}}
 {2} \)}{\(0\)}{\(1\)}{\(-\frac{\sqrt{3}} 
 {3} \)}{\(-\frac{\sqrt{2}} 
 {2} \)}}
\LANGes{
\Question{
\(\cos \left (\pi \right ) = ?\)}
{\(- 1\)}{\(\sqrt{3}\)}{\(\frac{\sqrt{3}}
 {3} \)}{\(-\sqrt{3}\)}{\(-\frac{\sqrt{3}} 
 {3} \)}{\(0\)}}
\End

\Data\ID{9000032108}
\Updated{2020-07-15 03:07:34}
\Level{A}
\SubArea{Sine, cosine, tangent and cotangent}
\IMG{000321_08-figure0.pdf}
\LANGen{
\Question{
Evaluate
\(\cos \left (\frac{-3\pi }
 {2} \right )\).}
{\(0\)}{\(-\frac{\sqrt{2}} 
 {2} \)}{\(- 1\)}{\(1\)}{\(-\frac{\sqrt{3}} 
 {3} \)}{\(\sqrt{3}\)}}
\LANGpl{
\Question{
Oblicz
\(\cos \left (\frac{-3\pi }
 {2} \right )\).}
{\(0\)}{\(-\frac{\sqrt{2}} 
 {2} \)}{\(- 1\)}{\(1\)}{\(-\frac{\sqrt{3}} 
 {3} \)}{\(\sqrt{3}\)}}
\LANGcs{
\Question{
\(\cos \left (\frac{-3\pi }
 {2} \right ) =?\)}
{\(0\)}{\(1\)}{\(-\frac{\sqrt{3}} 
 {3} \)}{\(\frac{\sqrt{3}}
 {3} \)}{\(-\sqrt{3}\)}{\(\sqrt{3}\)}}
\LANGsk{
\Question{
\(\cos \left (\frac{-3\pi }
 {2} \right ) =?\)}
{\(0\)}{\(1\)}{\(-\frac{\sqrt{3}} 
 {3} \)}{\(\frac{\sqrt{3}}
 {3} \)}{\(-\sqrt{3}\)}{\(\sqrt{3}\)}}
\LANGes{
\Question{
\(\cos \left (\frac{-3\pi }
 {2} \right) = ?\)}
{\(0\)}{\(-\frac{\sqrt{2}} 
 {2} \)}{\(- 1\)}{\(1\)}{\(-\frac{\sqrt{3}} 
 {3} \)}{\(\sqrt{3}\)}}
\End

\Data\ID{9000032109}
\Updated{2020-07-15 03:07:34}
\Level{A}
\SubArea{Sine, cosine, tangent and cotangent}
\IMG{000321_09-figure0.pdf}
\LANGen{
\Question{
Evaluate
\(\sin \left ( \frac{\pi }{4}\right )\).}
{\(\frac{\sqrt{2}}
 {2} \)}{\(-\frac{\sqrt{2}} 
 {2} \)}{\(0\)}{\(-\frac{\sqrt{3}} 
 {3} \)}{\(\frac{\sqrt{3}}
 {3} \)}{\(1\)}}
\LANGpl{
\Question{
Oblicz
\(\sin \left ( \frac{\pi }{4}\right )\).}
{\(\frac{\sqrt{2}}
 {2} \)}{\(-\frac{\sqrt{2}} 
 {2} \)}{\(0\)}{\(-\frac{\sqrt{3}} 
 {3} \)}{\(\frac{\sqrt{3}}
 {3} \)}{\(1\)}}
\LANGcs{
\Question{
\(\sin \left ( \frac{\pi }{4}\right ) =?\)}
{\(\frac{\sqrt{2}}
 {2} \)}{\(\sqrt{3}\)}{\(-\frac{\sqrt{2}} 
 {2} \)}{\(1\)}{\(- 1\)}{\(0\)}}
\LANGsk{
\Question{
\(\sin \left ( \frac{\pi }{4}\right ) =?\)}
{\(\frac{\sqrt{2}}
 {2} \)}{\(\sqrt{3}\)}{\(-\frac{\sqrt{2}} 
 {2} \)}{\(1\)}{\(- 1\)}{\(0\)}}
\LANGes{
\Question{
\(\sin \left ( \frac{\pi }{4}\right) = ?\)}
{\(\frac{\sqrt{2}}
 {2} \)}{\(-\frac{\sqrt{2}} 
 {2} \)}{\(0\)}{\(-\frac{\sqrt{3}} 
 {3} \)}{\(\frac{\sqrt{3}}
 {3} \)}{\(1\)}}
\End

\Data\ID{9000032110}
\Updated{2020-07-15 03:07:34}
\Level{A}
\SubArea{Sine, cosine, tangent and cotangent}
\IMG{000321_10-figure0.pdf}
\LANGen{
\Question{
Evaluate
\(\cos \left ( \frac{\pi }{4}\right )\).}
{\(\frac{\sqrt{2}}
 {2} \)}{\(0\)}{\(-\sqrt{3}\)}{\(-\frac{\sqrt{2}} 
 {2} \)}{\(- 1\)}{\(\sqrt{3}\)}}
\LANGpl{
\Question{
Oblicz
\(\cos \left ( \frac{\pi }{4}\right )\).}
{\(\frac{\sqrt{2}}
 {2} \)}{\(0\)}{\(-\sqrt{3}\)}{\(-\frac{\sqrt{2}} 
 {2} \)}{\(- 1\)}{\(\sqrt{3}\)}}
\LANGcs{
\Question{
\(\cos \left ( \frac{\pi }{4}\right ) =?\)}
{\(\frac{\sqrt{2}}
 {2} \)}{\(1\)}{\(\sqrt{3}\)}{\(0\)}{\(-\frac{\sqrt{2}} 
 {2} \)}{\(-\sqrt{3}\)}}
\LANGsk{
\Question{
\(\cos \left ( \frac{\pi }{4}\right ) =?\)}
{\(\frac{\sqrt{2}}
 {2} \)}{\(1\)}{\(\sqrt{3}\)}{\(0\)}{\(-\frac{\sqrt{2}} 
 {2} \)}{\(-\sqrt{3}\)}}
\LANGes{
\Question{
\(\cos \left ( \frac{\pi }{4}\right ) = ?\)}
{\(\frac{\sqrt{2}}
 {2} \)}{\(0\)}{\(-\sqrt{3}\)}{\(-\frac{\sqrt{2}} 
 {2} \)}{\(- 1\)}{\(\sqrt{3}\)}}
\End

\Data\ID{9000032111}
\Updated{2020-07-15 03:07:34}
\Level{A}
\SubArea{Sine, cosine, tangent and cotangent}
\IMG{000321_11-figure0.pdf}
\LANGen{
\Question{
Evaluate
\(\sin \left ( \frac{\pi }{3}\right )\).}
{\(\frac{\sqrt{3}}
 {2} \)}{\(-\frac{\sqrt{3}} 
 {2} \)}{\(\frac{\sqrt{2}}
 {2} \)}{\(0\)}{\(\sqrt{3}\)}{\(\frac{1}
{2}\)}}
\LANGpl{
\Question{
Oblicz
\(\sin \left ( \frac{\pi }{3}\right )\).}
{\(\frac{\sqrt{3}}
 {2} \)}{\(-\frac{\sqrt{3}} 
 {2} \)}{\(\frac{\sqrt{2}}
 {2} \)}{\(0\)}{\(\sqrt{3}\)}{\(\frac{1}
{2}\)}}
\LANGcs{
\Question{
\(\sin \left ( \frac{\pi }{3}\right ) =?\)}
{\(\frac{\sqrt{3}}
 {2} \)}{\(-\frac{\sqrt{3}} 
 {2} \)}{\(\frac{\sqrt{2}}
 {2} \)}{\(\sqrt{3}\)}{\(-\frac{1} 
{2}\)}{\(\frac{1}
{2}\)}}
\LANGsk{
\Question{
\(\sin \left ( \frac{\pi }{3}\right ) =?\)}
{\(\frac{\sqrt{3}}
 {2} \)}{\(-\frac{\sqrt{3}} 
 {2} \)}{\(\frac{\sqrt{2}}
 {2} \)}{\(\sqrt{3}\)}{\(-\frac{1} 
{2}\)}{\(\frac{1}
{2}\)}}
\LANGes{
\Question{
\(\sin \left ( \frac{\pi }{3}\right ) = ?\)}
{\(\frac{\sqrt{3}}
 {2} \)}{\(-\frac{\sqrt{3}} 
 {2} \)}{\(\frac{\sqrt{2}}
 {2} \)}{\(0\)}{\(\sqrt{3}\)}{\(\frac{1}
{2}\)}}
\End

\Data\ID{9000032112}
\Updated{2020-07-15 03:07:34}
\Level{A}
\SubArea{Sine, cosine, tangent and cotangent}
\IMG{000321_12-figure0.pdf}
\LANGen{
\Question{
Evaluate
\(\cos \left ( \frac{\pi }{3}\right )\).}
{\(\frac{1}
{2}\)}{\(\frac{\sqrt{3}}
 {2} \)}{\(-\frac{\sqrt{3}} 
 {2} \)}{\(\sqrt{3}\)}{\(-\frac{\sqrt{2}} 
 {2} \)}{\(\frac{\sqrt{2}}
 {2} \)}}
\LANGpl{
\Question{
Oblicz
\(\cos \left ( \frac{\pi }{3}\right )\).}
{\(\frac{1}
{2}\)}{\(\frac{\sqrt{3}}
 {2} \)}{\(-\frac{\sqrt{3}} 
 {2} \)}{\(\sqrt{3}\)}{\(-\frac{\sqrt{2}} 
 {2} \)}{\(\frac{\sqrt{2}}
 {2} \)}}
\LANGcs{
\Question{
\(\cos \left ( \frac{\pi }{3}\right ) =?\)}
{\(\frac{1}
{2}\)}{\(\sqrt{3}\)}{\(\frac{\sqrt{3}}
 {2} \)}{\(-\frac{\sqrt{3}} 
 {2} \)}{\(\frac{\sqrt{2}}
 {2} \)}{\(-\frac{\sqrt{2}} 
 {2} \)}}
\LANGsk{
\Question{
\(\cos \left ( \frac{\pi }{3}\right ) =?\)}
{\(\frac{1}
{2}\)}{\(\sqrt{3}\)}{\(\frac{\sqrt{3}}
 {2} \)}{\(-\frac{\sqrt{3}} 
 {2} \)}{\(\frac{\sqrt{2}}
 {2} \)}{\(-\frac{\sqrt{2}} 
 {2} \)}}
\LANGes{
\Question{
\(\cos \left ( \frac{\pi }{3}\right ) = ?\)}
{\(\frac{1}
{2}\)}{\(\frac{\sqrt{3}}
 {2} \)}{\(-\frac{\sqrt{3}} 
 {2} \)}{\(\sqrt{3}\)}{\(-\frac{\sqrt{2}} 
 {2} \)}{\(\frac{\sqrt{2}}
 {2} \)}}
\End

\Data\ID{9000032113}
\Updated{2020-07-15 03:07:34}
\Level{A}
\SubArea{Sine, cosine, tangent and cotangent}
\IMG{000321_13-figure0.pdf}
\LANGen{
\Question{
Evaluate
\(\sin \left ( \frac{\pi }{6}\right )\).}
{\(\frac{1}
{2}\)}{\(\frac{\sqrt{2}}
 {2} \)}{\(-\frac{1} 
{2}\)}{\(\frac{\sqrt{3}}
 {2} \)}{\(\sqrt{3}\)}{\(0\)}}
\LANGpl{
\Question{
Oblicz
\(\sin \left ( \frac{\pi }{6}\right )\).}
{\(\frac{1}
{2}\)}{\(\frac{\sqrt{2}}
 {2} \)}{\(-\frac{1} 
{2}\)}{\(\frac{\sqrt{3}}
 {2} \)}{\(\sqrt{3}\)}{\(0\)}}
\LANGcs{
\Question{
\(\sin \left ( \frac{\pi }{6}\right ) =?\)}
{\(\frac{1}
{2}\)}{\(-\sqrt{3}\)}{\(\frac{\sqrt{2}}
 {2} \)}{\(-\frac{1} 
{2}\)}{\(\frac{\sqrt{3}}
 {2} \)}{\(-\frac{\sqrt{2}} 
 {2} \)}}
\LANGsk{
\Question{
\(\sin \left ( \frac{\pi }{6}\right ) =?\)}
{\(\frac{1}
{2}\)}{\(-\sqrt{3}\)}{\(\frac{\sqrt{2}}
 {2} \)}{\(-\frac{1} 
{2}\)}{\(\frac{\sqrt{3}}
 {2} \)}{\(-\frac{\sqrt{2}} 
 {2} \)}}
\LANGes{
\Question{
\(\sin \left ( \frac{\pi }{6}\right ) = ?\)}
{\(\frac{1}
{2}\)}{\(\frac{\sqrt{2}}
 {2} \)}{\(-\frac{1} 
{2}\)}{\(\frac{\sqrt{3}}
 {2} \)}{\(\sqrt{3}\)}{\(0\)}}
\End

\Data\ID{9000032114}
\Updated{2020-07-15 03:07:34}
\Level{A}
\SubArea{Sine, cosine, tangent and cotangent}
\IMG{000321_14-figure0.pdf}
\LANGen{
\Question{
Evaluate
\(\cos \left ( \frac{\pi }{6}\right )\).}
{\(\frac{\sqrt{3}}
 {2} \)}{\(\frac{\sqrt{2}}
 {2} \)}{\(-\frac{1} 
{2}\)}{\(-\frac{\sqrt{3}} 
 {2} \)}{\(0\)}{\(-\frac{\sqrt{2}} 
 {2} \)}}
\LANGpl{
\Question{
Oblicz
\(\cos \left ( \frac{\pi }{6}\right )\).}
{\(\frac{\sqrt{3}}
 {2} \)}{\(\frac{\sqrt{2}}
 {2} \)}{\(-\frac{1} 
{2}\)}{\(-\frac{\sqrt{3}} 
 {2} \)}{\(0\)}{\(-\frac{\sqrt{2}} 
 {2} \)}}
\LANGcs{
\Question{
\(\cos \left ( \frac{\pi }{6}\right ) =?\)}
{\(\frac{\sqrt{3}}
 {2} \)}{\(-\frac{\sqrt{3}} 
 {2} \)}{\(\frac{\sqrt{2}}
 {2} \)}{\(\frac{1}
{2}\)}{\(0\)}{\(-\frac{1} 
{2}\)}}
\LANGsk{
\Question{
\(\cos \left ( \frac{\pi }{6}\right ) =?\)}
{\(\frac{\sqrt{3}}
 {2} \)}{\(-\frac{\sqrt{3}} 
 {2} \)}{\(\frac{\sqrt{2}}
 {2} \)}{\(\frac{1}
{2}\)}{\(0\)}{\(-\frac{1} 
{2}\)}}
\LANGes{
\Question{
\(\cos \left ( \frac{\pi }{6}\right) = ?\)}
{\(\frac{\sqrt{3}}
 {2} \)}{\(\frac{\sqrt{2}}
 {2} \)}{\(-\frac{1} 
{2}\)}{\(-\frac{\sqrt{3}} 
 {2} \)}{\(0\)}{\(-\frac{\sqrt{2}} 
 {2} \)}}
\End

\Data\ID{1003048506}
\Updated{2021-04-24 03:04:43}
\Level{B}
\SubArea{Sine, cosine, tangent and cotangent}
\LANGen{
\Question{
Identify which of the following functions has the smallest period.}
{\( f(x)=(\cos(2x) )^2 \)}{\( h(x)=\sin\bigl(\frac{x}{2}\bigr) \)}{\( m(x)=\mathrm{tg}\,\bigl(\frac{x}{2}\bigr) \)}{\( g(x)=(\mathrm{cotg}\, x)^2 \)}{}{}}
\LANGpl{
\Question{
Która z podanych funkcji ma najmniejszy okres?}
{\( f(x)=(\cos(2x) )^2 \)}{\( h(x)=\sin\bigl(\frac{x}{2}\bigr) \)}{\( m(x)=\mathrm{tg}\,\bigl(\frac{x}{2}\bigr) \)}{\( g(x)=(\mathrm{cotg}\, x)^2 \)}{}{}}
\LANGcs{
\Question{
Která z následujících funkcí má nejmenší periodu?}
{\( f(x)=(\cos(2x) )^2 \)}{\( h(x)=\sin\bigl(\frac{x}{2}\bigr) \)}{\( m(x)=\mathrm{tg}\,\bigl(\frac{x}{2}\bigr) \)}{\( g(x)=(\mathrm{cotg}\, x)^2 \)}{}{}}
\LANGsk{
\Question{
Ktorá z následujúcich funkcií má najmenšiu periódu?}
{\( f(x)=(\cos(2x) )^2 \)}{\( h(x)=\sin\bigl(\frac{x}{2}\bigr) \)}{\( m(x)=\mathrm{tg}\,\bigl(\frac{x}{2}\bigr) \)}{\( g(x)=(\mathrm{cotg}\, x)^2 \)}{}{}}
\LANGes{
\Question{
Identifica cuál de las siguientes funciones tiene menor período.}
{\( f(x)=(\cos(2x) )^2 \)}{\( h(x)=\sin\bigl(\frac{x}{2}\bigr) \)}{\( m(x)=\mathrm{tg}\,\bigl(\frac{x}{2}\bigr) \)}{\( g(x)=(\mathrm{cotg}\, x)^2 \)}{}{}}
\End

\Data\ID{1003076501}
\Updated{2020-07-15 03:07:12}
\Level{B}
\SubArea{Sine, cosine, tangent and cotangent}
\LANGen{
\Question{
To which quadrant does the angle \( \alpha \) belong if \( \sin\alpha=0.8 \) and \( \cos\alpha < 0 \)?}
{II.}{I.}{III.}{IV.}{}{}}
\LANGpl{
\Question{
Do której ćwiartki należy kąt \( \alpha \) jeżeli \( \sin\alpha=0{,}8 \) i \( \cos\alpha < 0 \)?}
{II.}{I.}{III.}{IV.}{}{}}
\LANGcs{
\Question{
Do kterého kvadrantu patří úhel \( \alpha \), jestliže \( \sin\alpha=0{,}8 \) a \( \cos\alpha < 0 \)?}
{II.}{I.}{III.}{IV.}{}{}}
\LANGsk{
\Question{
Do ktorého kvadrantu patrí uhol  \( \alpha \)   ak \( \sin\alpha=0{,}8 \) a \( \cos\alpha < 0 \)?}
{II.}{I.}{III.}{IV.}{}{}}
\LANGes{
\Question{
¿A qué cuadrante pertenece el ángulo \(\ alpha\) si \( \sin\alpha=0.8 \) y \( \cos\alpha < 0 \)?}
{II.}{I.}{III.}{IV.}{}{}}
\End

\Data\ID{1003076502}
\Updated{2020-07-15 03:07:12}
\Level{B}
\SubArea{Sine, cosine, tangent and cotangent}
\LANGen{
\Question{
To which quadrant does the angle \( \alpha \) belong if \( \sin\alpha< 0  \) and \( \cos\alpha < 0 \)?}
{III.}{I.}{II.}{IV.}{}{}}
\LANGpl{
\Question{
Do której ćwiartki należy kąt \( \alpha \) jeżeli \( \sin\alpha < 0\) i \( \cos\alpha < 0 \)?}
{III.}{I.}{II.}{IV.}{}{}}
\LANGcs{
\Question{
Do kterého kvadrantu patří úhel \( \alpha \), jestliže \( \sin\alpha < 0  \) a \( \cos\alpha < 0 \)?}
{III.}{I.}{II.}{IV.}{}{}}
\LANGsk{
\Question{
Do ktorého kvadrantu patrí uhol  \( \alpha \)   ak \( \sin\alpha < 0  \) a \( \cos\alpha < 0 \)?}
{III.}{I.}{II.}{IV.}{}{}}
\LANGes{
\Question{
¿A qué cuadrante pertenece el ángulo  \( \alpha \) si \( \sin\alpha< 0  \) y \( \cos\alpha < 0 \)?}
{III.}{I.}{II.}{IV.}{}{}}
\End

\Data\ID{1003076503}
\Updated{2021-04-24 03:04:35}
\Level{B}
\SubArea{Sine, cosine, tangent and cotangent}
\LANGen{
\Question{
Choose the true statement that applies to each of the functions \( f(x)=\sin x \), \( g(x)=\cos x \), \( h(x)= \mathrm{tg}\,x \):}
{The function has infinitely many zero points.}{The function is odd.}{The function is bounded.}{The function is simple.}{}{}}
\LANGpl{
\Question{
Wybierz stwierdzenie dotyczące funkcji  \( f(x)=\sin x \), \( g(x)=\cos x \), \( h(x)= \mathrm{tg}\,x \):}
{Funkcja ma nieskończenie wiele miejsc zerowych.}{Funkcja jest nieparzysta.}{Funkcja jest ograniczona.}{Funkcja jest prosta.}{}{}}
\LANGcs{
\Question{
Vyberte pravdivé tvrzení, které platí pro každou z funkcí \( f(x)=\sin x \), \( g(x)=\cos x \), \( h(x)= \mathrm{tg}\,x \):}
{Funkce má nekonečně mnoho nulových bodů.}{Funkce je lichá.}{Funkce je ohraničená.}{Funkce je prostá.}{}{}}
\LANGsk{
\Question{
Vyberte pravdivé tvrdenie, ktoré platí pre každú z funkcií  \( f(x)=\sin x \), \( g(x)=\cos x \), \( h(x)= \mathrm{tg}\,x \):}
{Funkcia má nekonečne veľa nulových bodov.}{Funkcia je nepárna.}{Funkcia je ohraničená.}{Funkcia je prostá.}{}{}}
\LANGes{
\Question{
Elige la proposición verdadera que es válida para cada  cada una de las funciones \( f(x)=\sin x \), \( g(x)=\cos x \), \( h(x)= \mathrm{tg}\,x \):}
{La función tiene infinitos puntos nulos.}{La función es impar.}{La función está acotada.}{La función es simple.}{}{}}
\End

\Data\ID{1003076504}
\Updated{2020-07-15 03:07:12}
\Level{B}
\SubArea{Sine, cosine, tangent and cotangent}
\LANGen{
\Question{
Choose the false statement:}
{The function \( f(x)= \mathrm{tg}\,x \) is even.}{The function \( f(x)=\mathrm{cotg}\,x \) is decreasing on the interval \( (0;\pi) \).}{The function \( f(x)=\sin x \) is bounded on its domain.}{The values of function \( f(x)=\cos x \)  always lie between \( -1 \) and \( 1 \).}{}{}}
\LANGpl{
\Question{
Wybierz fałszywe stwierdzenie:}
{Funkcja \( f(x)= \mathrm{tg}\,x \) jest parzysta.}{Funkcja  \( f(x)=\mathrm{cotg}\,x \) jest malejąca w przedziale  \( (0;\pi) \).}{Funkcja  \( f(x)=\sin x \) jest ograniczona w dziedzinie.}{Wartości funkcji \( f(x)=\cos x \)  zawsze leżą pomiędzy  \( -1 \) i \( 1 \).}{}{}}
\LANGcs{
\Question{
Vyberte nepravdivé tvrzení:}
{Funkce \( f(x)= \mathrm{tg}\,x \) je sudá.}{Funkce \( f(x)=\mathrm{cotg}\,x \)  je v intervalu \( (0;\pi) \) klesající.}{Funkce \( f(x)=\sin x\) je ohraničená v celém definičním oboru.}{Funkce \( f(x)=\cos x \)  nabývá funkční hodnoty od  \( -1 \) po  \( 1 \).}{}{}}
\LANGsk{
\Question{
Vyberte nepravdivé tvrdenie:}
{Funkcia \( f(x)= \mathrm{tg}\,x \) je párna.}{Funkcia \( f(x)=\mathrm{cotg}\,x \)  je na intervale \( (0;\pi) \) klesajúca.}{Funkcia \( f(x)=\sin x\) je ohraničená v celom definičnom obore.}{Funkcia  \( f(x)=\cos x \)  nadobúda funkčné hodnoty od  \( -1 \) po  \( 1 \).}{}{}}
\LANGes{
\Question{
Elige una proposición falsa:}
{La función \( f(x)= \mathrm{tg}\,x \) es par.}{La función \( f(x)=\mathrm{cotg}\,x \) es decreciente en el intervalo \( (0;\pi) \).}{La función \( f(x)=\sin x \) está acotada en su dominio.}{Los valores de la función \( f(x)=\cos x \)  se encuentran siempre entre \( -1 \) y  \( 1 \).}{}{}}
\End

\Data\ID{1003076505}
\Updated{2021-04-24 04:04:27}
\Level{B}
\SubArea{Sine, cosine, tangent and cotangent}
\LANGen{
\Question{
Choose the false statement:}
{\( \cos190^{\circ} > \cos240^{\circ} \)}{\( \sin140^{\circ} >\sin190^{\circ} \)}{\( \sin15^{\circ}>\sin210^{\circ} \)}{\( \cos305^{\circ}>\cos300^{\circ} \)}{}{}}
\LANGpl{
\Question{
Wybierz fałszywe wyrażenie:}
{\( \cos190^{\circ} > \cos240^{\circ} \)}{\( \sin140^{\circ} >\sin190^{\circ} \)}{\( \sin15^{\circ}>\sin210^{\circ} \)}{\( \cos305^{\circ}>\cos300^{\circ} \)}{}{}}
\LANGcs{
\Question{
Vyberte nepravdivé tvrzení:}
{\( \cos190^{\circ} > \cos240^{\circ} \)}{\( \sin140^{\circ} >\sin190^{\circ} \)}{\( \sin15^{\circ}>\sin210^{\circ} \)}{\( \cos305^{\circ}>\cos300^{\circ} \)}{}{}}
\LANGsk{
\Question{
Vyberte nepravdivé tvrdenie:}
{\( \cos190^{\circ} > \cos240^{\circ} \)}{\( \sin140^{\circ} >\sin190^{\circ} \)}{\( \sin15^{\circ}>\sin210^{\circ} \)}{\( \cos305^{\circ}>\cos300^{\circ} \)}{}{}}
\LANGes{
\Question{
Elige la desigualdad falsa:}
{\( \cos190^{\circ} > \cos240^{\circ} \)}{\( \sin140^{\circ} >\sin190^{\circ} \)}{\( \sin15^{\circ}>\sin210^{\circ} \)}{\( \cos305^{\circ}>\cos300^{\circ} \)}{}{}}
\End

\Data\ID{1003076506}
\Updated{2020-07-15 03:07:12}
\Level{B}
\SubArea{Sine, cosine, tangent and cotangent}
\LANGen{
\Question{
Find the smallest period of the function \( f(x)=\mathrm{tg}\,4x \):}
{\( \frac{\pi}4 \)}{\( 4\pi \)}{\( \pi \)}{\( 2\pi \)}{}{}}
\LANGpl{
\Question{
Wskaż  najmniejszy okres funkcji \( f(x)=\mathrm{tg}\,4x \):}
{\( \frac{\pi}4 \)}{\( 4\pi \)}{\( \pi \)}{\( 2\pi \)}{}{}}
\LANGcs{
\Question{
Najdi nejmenší periodu funkce \( f(x)=\mathrm{tg}\,4x \):}
{\( \frac{\pi}4 \)}{\( 4\pi \)}{\( \pi \)}{\( 2\pi \)}{}{}}
\LANGsk{
\Question{
Nájdite najmenšiu periódu funkcie  \( f(x)=\mathrm{tg}\,4x \):}
{\( \frac{\pi}4 \)}{\( 4\pi \)}{\( \pi \)}{\( 2\pi \)}{}{}}
\LANGes{
\Question{
Determina el período más pequeño de la función \( f(x)=\mathrm{tg}\,4x \):}
{\( \frac{\pi}4 \)}{\( 4\pi \)}{\( \pi \)}{\( 2\pi \)}{}{}}
\End

\Data\ID{1003076507}
\Updated{2021-04-24 04:04:31}
\Level{B}
\SubArea{Sine, cosine, tangent and cotangent}
\LANGen{
\Question{
For which \( x \in[0;\frac{\pi}2] \) is \( \sin x = \cos x \)?}
{\( \frac{\pi}4 \)}{\( 0 \)}{\( \frac{\pi}2 \)}{\( \frac{\pi}3 \)}{}{}}
\LANGpl{
\Question{
Dla jakich \( x \in\langle0;\frac{\pi}2\rangle \) jest  \( \sin x = \cos x \)?}
{\( \frac{\pi}4 \)}{\( 0 \)}{\( \frac{\pi}2 \)}{\( \frac{\pi}3 \)}{}{}}
\LANGcs{
\Question{
Pro které \( x \in\langle0;\frac{\pi}2\rangle \) platí \( \sin x = \cos x \)?}
{\( \frac{\pi}4 \)}{\( 0 \)}{\( \frac{\pi}2 \)}{\( \frac{\pi}3 \)}{}{}}
\LANGsk{
\Question{
Pre ktoré \( x \in\langle0;\frac{\pi}2\rangle \) je \( \sin x = \cos x \)?}
{\( \frac{\pi}4 \)}{\( 0 \)}{\( \frac{\pi}2 \)}{\( \frac{\pi}3 \)}{}{}}
\LANGes{
\Question{
Para qué valores de \( x \in[0;\frac{\pi}2] \) es \( \sin x = \cos x \)?}
{\( \frac{\pi}4 \)}{\( 0 \)}{\( \frac{\pi}2 \)}{\( \frac{\pi}3 \)}{}{}}
\End

\Data\ID{1003076508}
\Updated{2021-04-24 04:04:53}
\Level{B}
\SubArea{Sine, cosine, tangent and cotangent}
\LANGen{
\Question{
For which \( \alpha\in[0;90^{\circ} ] \) is \( \mathrm{tg}\,\alpha = \mathrm{cotg}\,\alpha \)?}
{\( 45^{\circ} \)}{\( 60^{\circ} \)}{\( 35^{\circ} \)}{\( 30^{\circ} \)}{}{}}
\LANGpl{
\Question{
Dla jakich \( \alpha\in\langle0;90^{\circ} \rangle \) jest \( \mathrm{tg}\,\alpha = \mathrm{cotg}\,\alpha \)?}
{\( 45^{\circ} \)}{\( 60^{\circ} \)}{\( 35^{\circ} \)}{\( 30^{\circ} \)}{}{}}
\LANGcs{
\Question{
Pro které \( \alpha\in\langle0;90^{\circ} \rangle \) platí \( \mathrm{tg}\,\alpha = \mathrm{cotg}\,\alpha \)?}
{\( 45^{\circ} \)}{\( 60^{\circ} \)}{\( 35^{\circ} \)}{\( 30^{\circ} \)}{}{}}
\LANGsk{
\Question{
Pre ktoré  \( \alpha\in\langle0;90^{\circ} \rangle \) je \( \mathrm{tg}\,\alpha = \mathrm{cotg}\,\alpha \)?}
{\( 45^{\circ} \)}{\( 60^{\circ} \)}{\( 35^{\circ} \)}{\( 30^{\circ} \)}{}{}}
\LANGes{
\Question{
Para qué valores de \( \alpha\in[0;90^{\circ} ] \) es \( \mathrm{tg}\,\alpha = \mathrm{cotg}\,\alpha \)?}
{\( 45^{\circ} \)}{\( 60^{\circ} \)}{\( 35^{\circ} \)}{\( 30^{\circ} \)}{}{}}
\End

\Data\ID{1003076509}
\Updated{2020-07-15 03:07:12}
\Level{B}
\SubArea{Sine, cosine, tangent and cotangent}
\LANGen{
\Question{
At which point \( x\in(2\pi; 4\pi) \) has the function \( f(x)=\cos x \) a minimum?}
{\(3\pi\)}{\(2 \pi \)}{\( 4\pi \)}{\( 3.5\pi \)}{}{}}
\LANGpl{
\Question{
W którym punkcie \( x\in(2\pi; 4\pi) \)  funkcja \( f(x)=\cos x \) ma minimum?}
{\(3\pi\)}{\(2 \pi \)}{\( 4\pi \)}{\( 3{,}5\pi \)}{}{}}
\LANGcs{
\Question{
Pro které \( x\in(2\pi; 4\pi) \) nabývá funkce \( f(x)=\cos x \) minimum?}
{\(3\pi\)}{\(2 \pi \)}{\( 4\pi \)}{\( 3{,}5\pi \)}{}{}}
\LANGsk{
\Question{
Pre ktoré \( x\in(2\pi; 4\pi) \)  nadobúda funkcia  \( f(x)=\cos x \) minimum?}
{\(3\pi\)}{\(2 \pi \)}{\( 4\pi \)}{\( 3{,}5\pi \)}{}{}}
\LANGes{
\Question{
En qué punto  \( x\in(2\pi; 4\pi) \) tiene la función \( f(x)=\cos x \) un mínimo?}
{\(3\pi\)}{\(2 \pi \)}{\( 4\pi \)}{\( 3.5\pi \)}{}{}}
\End

\Data\ID{1003076510}
\Updated{2020-07-15 03:07:12}
\Level{B}
\SubArea{Sine, cosine, tangent and cotangent}
\LANGen{
\Question{
If the angle \( \alpha\in[0;90^{\circ}] \) and \( \sin\alpha = 0.5 \), then \( \cos \alpha \) is equal to:}
{\( \frac{\sqrt3}2 \)}{\( -\frac{\sqrt{3}}2 \)}{\( -\frac12 \)}{\( \frac12 \)}{}{}}
\LANGpl{
\Question{
Jeśli kąt \( \alpha\in\langle0;90^{\circ}\rangle \) i \( \sin\alpha = 0{,}5 \), wtedy \( \cos \alpha \) jest równy:}
{\( \frac{\sqrt3}2 \)}{\( -\frac{\sqrt{3}}2 \)}{\( -\frac12 \)}{\( \frac12 \)}{}{}}
\LANGcs{
\Question{
Jestliže úhel \( \alpha\in\langle0;90^{\circ}\rangle \) a \( \sin\alpha = 0{,}5 \), pak \( \cos \alpha \)  se rovná:}
{\( \frac{\sqrt3}2 \)}{\( -\frac{\sqrt{3}}2 \)}{\( -\frac12 \)}{\( \frac12 \)}{}{}}
\LANGsk{
\Question{
Ak uhol   \( \alpha\in\langle0;90^{\circ}\rangle \) a \( \sin\alpha = 0{,}5 \), tak \( \cos \alpha \)  sa rovná:}
{\( \frac{\sqrt3}2 \)}{\( -\frac{\sqrt{3}}2 \)}{\( -\frac12 \)}{\( \frac12 \)}{}{}}
\LANGes{
\Question{
Si el ángulo \( \alpha\in[0;90^{\circ}] \) y \( \sin\alpha = 0.5 \), entonces \( \cos \alpha \) es igual a:}
{\( \frac{\sqrt3}2 \)}{\( -\frac{\sqrt{3}}2 \)}{\( -\frac12 \)}{\( \frac12 \)}{}{}}
\End

\Data\ID{1003076511}
\Updated{2020-07-15 03:07:12}
\Level{B}
\SubArea{Sine, cosine, tangent and cotangent}
\LANGen{
\Question{
If \( \mathrm{tg}\,x = 1 \), then \( \mathrm{cotg}\,x \) equals:}
{\( 1 \)}{\( 0 \)}{\( -1 \)}{\( 0.5 \)}{}{}}
\LANGpl{
\Question{
Jeżeli  \( \mathrm{tg}\,x = 1 \), wtedy \( \mathrm{cotg}\,x \) jest równy:}
{\( 1 \)}{\( 0 \)}{\( -1 \)}{\( 0{,}5 \)}{}{}}
\LANGcs{
\Question{
Jestliže \( \mathrm{tg}\,x = 1 \), pak \( \mathrm{cotg}\,x \) se rovná:}
{\( 1 \)}{\( 0 \)}{\( -1 \)}{\( 0{,}5 \)}{}{}}
\LANGsk{
\Question{
Ak \( \mathrm{tg}\,x = 1 \), tak \( \mathrm{cotg}\,x \) sa rovná:}
{\( 1 \)}{\( 0 \)}{\( -1 \)}{\( 0{,}5 \)}{}{}}
\LANGes{
\Question{
Si \( \mathrm{tg}\,x = 1 \), entonces \( \mathrm{cotg}\,x \) equivale a:}
{\( 1 \)}{\( 0 \)}{\( -1 \)}{\( 0.5 \)}{}{}}
\End

\Data\ID{1003076512}
\Updated{2020-07-15 03:07:12}
\Level{B}
\SubArea{Sine, cosine, tangent and cotangent}
\LANGen{
\Question{
The value of \( \sin\left(-\frac{53\pi}6\right) \) is the same as the value of}
{\( \sin\frac{7\pi}6 \).}{\( \sin\frac{\pi}6 \).}{\( \sin\frac{5\pi}6 \).}{\( \sin\left( -\frac{11\pi}6 \right) \).}{}{}}
\LANGpl{
\Question{
Wartość  \( \sin\left(-\frac{53\pi}6\right) \) jest taka sama jak wartość}
{\( \sin\frac{7\pi}6 \).}{\( \sin\frac{\pi}6 \).}{\( \sin\frac{5\pi}6 \).}{\( \sin\left( -\frac{11\pi}6 \right) \).}{}{}}
\LANGcs{
\Question{
Hodnota \( \sin\left(-\frac{53\pi}6\right) \) je stejná jako hodnota}
{\( \sin\frac{7\pi}6 \).}{\( \sin\frac{\pi}6 \).}{\( \sin\frac{5\pi}6 \).}{\( \sin\left( -\frac{11\pi}6 \right) \).}{}{}}
\LANGsk{
\Question{
Výraz \( \sin\left(-\frac{53\pi}6\right) \)  má rovnakú hodnotu ako výraz}
{\( \sin\frac{7\pi}6 \).}{\( \sin\frac{\pi}6 \).}{\( \sin\frac{5\pi}6 \).}{\( \sin\left( -\frac{11\pi}6 \right) \).}{}{}}
\LANGes{
\Question{
El valor de \( \sin\left(-\frac{53\pi}6\right) \) es el mismo que el valor de}
{\( \sin\frac{7\pi}6 \).}{\( \sin\frac{\pi}6 \).}{\( \sin\frac{5\pi}6 \).}{\( \sin\left( -\frac{11\pi}6 \right) \).}{}{}}
\End

\Data\ID{1003076513}
\Updated{2020-07-15 03:07:12}
\Level{B}
\SubArea{Sine, cosine, tangent and cotangent}
\LANGen{
\Question{
If  two of the values of \( \sin\alpha \), \( \cos\alpha \), \( \mathrm{tg}\alpha\) and \( \mathrm{cotg}\alpha \) are negative, then \( \alpha \) belongs to the interval}
{\( \left(\pi;  \frac{3\pi}2 \right) \).}{\( \left(0; \frac{\pi}2 \right) \).}{\( \left(\frac{\pi}2; \pi\right) \).}{\( \left( \frac{3\pi}2; 2\pi \right) \).}{}{}}
\LANGpl{
\Question{
Jeżeli wartości \( \sin\alpha \), \( \cos\alpha \), \( \mathrm{tg}\alpha\) i \( \mathrm{cotg}\alpha \) są ujemne, wtedy \( \alpha \) należy do przedziału}
{\( \left(\pi;  \frac{3\pi}2 \right) \).}{\( \left(0; \frac{\pi}2 \right) \).}{\( \left(\frac{\pi}2; \pi\right) \).}{\( \left( \frac{3\pi}2; 2\pi \right) \).}{}{}}
\LANGcs{
\Question{
Pokud jsou dvě z hodnot \( \sin\alpha \), \( \cos\alpha \), \( \mathrm{tg}\alpha\) a \( \mathrm{cotg}\alpha \) záporné, pak \( \alpha \) náleží intervalu}
{\( \left(\pi;  \frac{3\pi}2 \right) \).}{\( \left(0; \frac{\pi}2 \right) \).}{\( \left(\frac{\pi}2; \pi\right) \).}{\( \left( \frac{3\pi}2; 2\pi \right) \).}{}{}}
\LANGsk{
\Question{
Ak  práve dve z hodnôt \( \sin\alpha \), \( \cos\alpha \), \( \mathrm{tg}\alpha\) and \( \mathrm{cotg}\alpha \)  sú záporné, potom \( \alpha \) patrí do  intervalu}
{\( \left(\pi;  \frac{3\pi}2 \right) \).}{\( \left(0; \frac{\pi}2 \right) \).}{\( \left(\frac{\pi}2; \pi\right) \).}{\( \left( \frac{3\pi}2; 2\pi \right) \).}{}{}}
\LANGes{
\Question{
Si dos de los valores de \( \sin\alpha \), \( \cos\alpha \), \( \mathrm{tg}\alpha\) y \( \mathrm{cotg}\alpha \) son negativos, entonces \( \alpha \) pertenece al intervalo}
{\( \left(\pi;  \frac{3\pi}2 \right) \).}{\( \left(0; \frac{\pi}2 \right) \).}{\( \left(\frac{\pi}2; \pi\right) \).}{\( \left( \frac{3\pi}2; 2\pi \right) \).}{}{}}
\End

\Data\ID{1003076514}
\Updated{2020-07-15 03:07:12}
\Level{B}
\SubArea{Sine, cosine, tangent and cotangent}
\LANGen{
\Question{
By simplifying the expression \( \cos\left( \frac{\pi}2  - x \right) \)  we get:}
{\( \sin x \)}{\( \cos x \)}{\( -\sin x \)}{\( -\cos x \)}{}{}}
\LANGpl{
\Question{
Uprość wyrażenie  \( \cos\left( \frac{\pi}2  - x \right) \).}
{\( \sin x \)}{\( \cos x \)}{\( -\sin x \)}{\( -\cos x \)}{}{}}
\LANGcs{
\Question{
Zjednodušením výrazu \( \cos\left( \frac{\pi}2  - x \right) \)  dostaneme:}
{\( \sin x \)}{\( \cos x \)}{\( -\sin x \)}{\( -\cos x \)}{}{}}
\LANGsk{
\Question{
Po úprave výrazu  \( \cos\left( \frac{\pi}2  - x \right) \) dostaneme:}
{\( \sin x \)}{\( \cos x \)}{\( -\sin x \)}{\( -\cos x \)}{}{}}
\LANGes{
\Question{
Al simplificar la expresión \( \cos\left( \frac{\pi}2  - x \right) \)  obtenemos:}
{\( \sin x \)}{\( \cos x \)}{\( -\sin x \)}{\( -\cos x \)}{}{}}
\End

\Data\ID{1003076515}
\Updated{2020-07-15 03:07:12}
\Level{B}
\SubArea{Sine, cosine, tangent and cotangent}
\LANGen{
\Question{
By simplifying the expression \( \sin\left(\frac{\pi}2  - x \right) \)  we get:}
{\( \cos x \)}{\( \sin x \)}{\( -\sin x \)}{\( -\cos x \)}{}{}}
\LANGpl{
\Question{
Uprość wyrażenie \( \sin\left(\frac{\pi}2  - x \right) \).}
{\( \cos x \)}{\( \sin x \)}{\( -\sin x \)}{\( -\cos x \)}{}{}}
\LANGcs{
\Question{
Zjednodušením výrazu \( \sin\left(\frac{\pi}2  - x \right) \)  dostaneme:}
{\( \cos x \)}{\( \sin x \)}{\( -\sin x \)}{\( -\cos x \)}{}{}}
\LANGsk{
\Question{
Po úprave výrazu  \( \sin\left(\frac{\pi}2  - x \right) \) dostaneme:}
{\( \cos x \)}{\( \sin x \)}{\( -\sin x \)}{\( -\cos x \)}{}{}}
\LANGes{
\Question{
Al simplificar la expresión \( \sin\left(\frac{\pi}2  - x \right) \) obtenemos:}
{\( \cos x \)}{\( \sin x \)}{\( -\sin x \)}{\( -\cos x \)}{}{}}
\End

\Data\ID{1003076601}
\Updated{2020-07-15 03:07:12}
\Level{B}
\SubArea{Sine, cosine, tangent and cotangent}
\LANGen{
\Question{
How many \( x \)-intercepts has the  graph of the function \( f(x)=\cos 2x \) on the interval \( [-\pi; 2\pi] \)?}
{\( 6 \)}{\( 4 \)}{\( 5 \)}{\( 3 \)}{}{}}
\LANGpl{
\Question{
Ile miejsc zerowych ma wykres funkcji  \( f(x)=\cos 2x \) w przedziale \( \langle-\pi; 2\pi\rangle \)?}
{\( 6 \)}{\( 4 \)}{\( 5 \)}{\( 3 \)}{}{}}
\LANGcs{
\Question{
Kolik průsečíků s osou  \( x \)  má graf funkce  \( f(x)=\cos 2x \) v intervalu \( \langle-\pi; 2\pi\rangle \)?}
{\( 6 \)}{\( 4 \)}{\( 5 \)}{\( 3 \)}{}{}}
\LANGsk{
\Question{
Koľko priesečníkov s osou  \( x \)  má graf funkcie  \( f(x)=\cos 2x \) na intervale \( \langle-\pi; 2\pi\rangle \)?}
{\( 6 \)}{\( 4 \)}{\( 5 \)}{\( 3 \)}{}{}}
\LANGes{
\Question{
¿Cuántas  intersecciones con el eje \( x \) tiene la gráfica de la función \( f(x)=\cos 2x \) en el intervalo \( [-\pi; 2\pi] \)?}
{\( 6 \)}{\( 4 \)}{\( 5 \)}{\( 3 \)}{}{}}
\End

\Data\ID{1003076602}
\Updated{2020-07-15 03:07:12}
\Level{B}
\SubArea{Sine, cosine, tangent and cotangent}
\LANGen{
\Question{
How many \( x \)-intercepts has the  graph of the function \( f(x)=\sin 3x \)  on the interval \( [-\pi; 3\pi] \)?}
{\( 13 \)}{\( 10 \)}{\( 14  \)}{\( 8 \)}{}{}}
\LANGpl{
\Question{
Ile miejsc zerowych ma wykres funkcji  \( f(x)=\sin 3x \)  w przedziale \( \langle-\pi; 3\pi\rangle \)?}
{\( 13 \)}{\( 10 \)}{\( 14  \)}{\( 8 \)}{}{}}
\LANGcs{
\Question{
Kolik průsečíků s osou \( x \) má  graf  funkce \( f(x)=\sin 3x \)  v intervalu \( \langle-\pi; 3\pi\rangle \)?}
{\( 13 \)}{\( 10 \)}{\( 14  \)}{\( 8 \)}{}{}}
\LANGsk{
\Question{
Koľko priesečníkov s osou \( x \) má  graf  funkcie \( f(x)=\sin 3x \)  na intervale \( \langle-\pi; 3\pi\rangle \)?}
{\( 13 \)}{\( 10 \)}{\( 14  \)}{\( 8 \)}{}{}}
\LANGes{
\Question{
¿Cuántas intersecciones con el eje \( x \) tiene la gráfica de la función \( f(x)=\sin 3x \)   en el intervalo \( [-\pi; 3\pi] \)?}
{\( 13 \)}{\( 10 \)}{\( 14  \)}{\( 8 \)}{}{}}
\End

\Data\ID{1003076603}
\Updated{2020-07-15 03:07:12}
\Level{B}
\SubArea{Sine, cosine, tangent and cotangent}
\LANGen{
\Question{
How many points of intersection has the graph  of the function \( f(x)=- 2\sin2x \) with the \( x \)-axis on the interval \( [-2\pi;  2\pi] \)?}
{\( 9 \)}{\( 8 \)}{\( 10  \)}{\( 11 \)}{}{}}
\LANGpl{
\Question{
Ile punktów przecięcia ma wykres funkcji  \( f(x)=- 2\sin2x \) z osią  \( x \) w przedziale \( \langle-2\pi;  2\pi\rangle \)?}
{\( 9 \)}{\( 8 \)}{\( 10  \)}{\( 11 \)}{}{}}
\LANGcs{
\Question{
Kolik společných bodů má graf funkce \( f(x)=- 2\sin2x \)  s osou \( x \)  v intervalu \( \langle-2\pi;  2\pi\rangle \)?}
{\( 9 \)}{\( 8 \)}{\( 10  \)}{\( 11 \)}{}{}}
\LANGsk{
\Question{
Koľko spoločných bodov má graf funkcie \( f(x)=- 2\sin2x \)  s  osou \( x \) na  intervale \( \langle-2\pi;  2\pi\rangle \)?}
{\( 9 \)}{\( 8 \)}{\( 10  \)}{\( 11 \)}{}{}}
\LANGes{
\Question{
¿Cuántas intersecciones con el eje \( x \) tiene la gráfica de la función \( f(x)=- 2\sin2x \) en el intervalo\( [-2\pi;  2\pi] \)?}
{\( 9 \)}{\( 8 \)}{\( 10  \)}{\( 11 \)}{}{}}
\End

\Data\ID{1003076604}
\Updated{2020-07-15 03:07:12}
\Level{B}
\SubArea{Sine, cosine, tangent and cotangent}
\LANGen{
\Question{
The graph of the function \( f(x)=-\sin (-x) \) is identical with the graph of the function:}
{\( g(x)=\cos\left(x -\frac{\pi}2 \right) \)}{\( g(x)=-\cos\left(x -\frac{\pi}2 \right) \)}{\( g(x)=\cos\left(x +\frac{\pi}2 \right) \)}{\( g(x)=-\sin x \)}{}{}}
\LANGpl{
\Question{
Wykres funkcji \( f(x)=-\sin (-x) \) jest identyczny z wykresem funkcji:}
{\( g(x)=\cos\left(x -\frac{\pi}2 \right) \)}{\( g(x)=-\cos\left(x -\frac{\pi}2 \right) \)}{\( g(x)=\cos\left(x +\frac{\pi}2 \right) \)}{\( g(x)=-\sin x \)}{}{}}
\LANGcs{
\Question{
Graf funkce \( f(x)=-\sin (-x) \) je shodný s grafem funkce:}
{\( g(x)=\cos\left(x -\frac{\pi}2 \right) \)}{\( g(x)=-\cos\left(x -\frac{\pi}2 \right) \)}{\( g(x)=\cos\left(x +\frac{\pi}2 \right) \)}{\( g(x)=-\sin x \)}{}{}}
\LANGsk{
\Question{
Graf funkcie \( f(x)=-\sin (-x) \)  je totožný s grafom funkcie:}
{\( g(x)=\cos\left(x -\frac{\pi}2 \right) \)}{\( g(x)=-\cos\left(x -\frac{\pi}2 \right) \)}{\( g(x)=\cos\left(x +\frac{\pi}2 \right) \)}{\( g(x)=-\sin x \)}{}{}}
\LANGes{
\Question{
La gráfica de la función \( f(x)=-\sin (-x) \) es idéntica a la gráfica de la función:}
{\( g(x)=\cos\left(x -\frac{\pi}2 \right) \)}{\( g(x)=-\cos\left(x -\frac{\pi}2 \right) \)}{\( g(x)=\cos\left(x +\frac{\pi}2 \right) \)}{\( g(x)=-\sin x \)}{}{}}
\End

\Data\ID{1003076605}
\Updated{2021-04-24 04:04:14}
\Level{B}
\SubArea{Sine, cosine, tangent and cotangent}
\LANGen{
\Question{
The graph of the function \( f(x)=\cos (-x) \) is identical with the graph  of the function:}
{\( g(x) = \cos x \)}{\( g(x) = \sin x \)}{\( g(x) = -\cos x \)}{\( g(x)=-\sin x \)}{}{}}
\LANGpl{
\Question{
Wykres funkcji  \( f(x)=\cos (-x) \) jest identyczny z wykresem funkcji:}
{\( g(x) = \cos x \)}{\( g(x) = \sin x \)}{\( g(x) = -\cos x \)}{\( g(x)=-\sin x \)}{}{}}
\LANGcs{
\Question{
Graf funkce \( f(x)=\cos (-x) \)  je shodný s grafem funkce:}
{\( g(x) = \cos x \)}{\( g(x) = \sin x \)}{\( g(x) = -\cos x \)}{\( g(x)=-\sin x \)}{}{}}
\LANGsk{
\Question{
Graf funkcie \( f(x)=\cos (-x) \)  je totožný s grafom funkcie:}
{\( g(x) = \cos x \)}{\( g(x) = \sin x \)}{\( g(x) = -\cos x \)}{\( g(x)=-\sin x \)}{}{}}
\LANGes{
\Question{
La gráfica de la función \( f(x)=\cos (-x) \) es idéntica a la gráfica de la función:}
{\( g(x) = \cos x \)}{\( g(x) = \sin x \)}{\( g(x) = -\cos x \)}{\( g(x)=-\sin x \)}{}{}}
\End

\Data\ID{1003076606}
\Updated{2020-07-15 03:07:12}
\Level{B}
\SubArea{Sine, cosine, tangent and cotangent}
\LANGen{
\Question{
The graph of the function \( g(x) = \cos x \) is identical with the graph of the function:}
{\( g(x) = \cos(-x) \)}{\( g(x) = \sin(- x ) \)}{\( g(x) = -\cos x \)}{\( g(x)= -\sin x \)}{}{}}
\LANGpl{
\Question{
Wykres funkcji  \( g(x) = \cos x \) jest identyczny z wykresem funkcji:}
{\( g(x) = \cos(-x) \)}{\( g(x) = \sin(- x ) \)}{\( g(x) = -\cos x \)}{\( g(x)= -\sin x \)}{}{}}
\LANGcs{
\Question{
Graf funkce \( g(x) = \cos x \) je shodný s grafem funkce:}
{\( g(x) = \cos(-x) \)}{\( g(x) = \sin(- x ) \)}{\( g(x) = -\cos x \)}{\( g(x)= -\sin x \)}{}{}}
\LANGsk{
\Question{
Graf funkcie \( g(x) = \cos x \) je totožný s grafom funkcie:}
{\( g(x) = \cos(-x) \)}{\( g(x) = \sin(- x ) \)}{\( g(x) = -\cos x \)}{\( g(x)= -\sin x \)}{}{}}
\LANGes{
\Question{
La gráfica de la función \( g(x) = \cos x \) es idéntica a la gráfica de la función:}
{\( g(x) = \cos(-x) \)}{\( g(x) = \sin(- x ) \)}{\( g(x) = -\cos x \)}{\( g(x)= -\sin x \)}{}{}}
\End

\Data\ID{1003076607}
\Updated{2020-07-15 03:07:12}
\Level{B}
\SubArea{Sine, cosine, tangent and cotangent}
\LANGen{
\Question{
If we reflect the graph of the function  \( f(x) = \cos x \) over the \( x \)-axis, we get the graph of the function:}
{\( g(x) = -\cos x \)}{\( g(x) =\sin x \)}{\( g(x) =\cos x \)}{\( g(x) = -\sin x \)}{}{}}
\LANGpl{
\Question{
Jeśli przekształcimy wykres funkcji  \( f(x) = \cos x \) względem osi \( x \),  otrzymamy wykres funkcji:}
{\( g(x) = -\cos x \)}{\( g(x) =\sin x \)}{\( g(x) =\cos x \)}{\( g(x) = -\sin x \)}{}{}}
\LANGcs{
\Question{
Graf funkce \( f(x) = \cos x \)  souměrně sdružíme podle osy \( x \). Dostaneme tak funkci, která je určená rovnicí:}
{\( g(x) = -\cos x \)}{\( g(x) =\sin x \)}{\( g(x) =\cos x \)}{\( g(x) = -\sin x \)}{}{}}
\LANGsk{
\Question{
Graf funkcie   \( f(x) = \cos x \)  súmerne združíme podľa osi  \( x \). Dostaneme tak funkciu, ktorá je určená rovnicou:}
{\( g(x) = -\cos x \)}{\( g(x) =\sin x \)}{\( g(x) =\cos x \)}{\( g(x) = -\sin x \)}{}{}}
\LANGes{
\Question{
Si reflejamos la gráfica de la función \(f (x) = \cos x\) sobre el eje \(x\), obtenemos la gráfica de la función:}
{\( g(x) = -\cos x \)}{\( g(x) =\sin x \)}{\( g(x) =\cos x \)}{\( g(x) = -\sin x \)}{}{}}
\End

\Data\ID{1003076706}
\Updated{2020-07-15 03:07:12}
\Level{B}
\SubArea{Sine, cosine, tangent and cotangent}
\LANGen{
\Question{
How many values of the angle \( \alpha\in\left(0^{\circ}; 90^{\circ}\right)\cup\left(90^{\circ}; 180^{\circ}\right) \)  satisfy the equation \( \mathrm{tg}\,\alpha = \mathrm{cotg}\,\alpha \)?}
{\( 2 \)}{\( 1 \)}{\( 0 \)}{\( 4 \)}{}{}}
\LANGpl{
\Question{
Jakie wartości kąta  \( \alpha\in\left(0^{\circ}; 90^{\circ}\right)\cup\left(90^{\circ}; 180^{\circ}\right) \)  spełniają równanie \( \mathrm{tg}\,\alpha = \mathrm{cotg}\,\alpha \)?}
{\( 2 \)}{\( 1 \)}{\( 0 \)}{\( 4 \)}{}{}}
\LANGcs{
\Question{
Kolik existuje úhlů \( \alpha\in\left(0^{\circ}; 90^{\circ}\right)\cup\left(90^{\circ}; 180^{\circ}\right) \), pro které \( \mathrm{tg}\,\alpha = \mathrm{cotg}\,\alpha \)?}
{\( 2 \)}{\( 1 \)}{\( 0 \)}{\( 4 \)}{}{}}
\LANGsk{
\Question{
Koľko je takých uhlov \( \alpha\in\left(0^{\circ}; 90^{\circ}\right)\cup\left(90^{\circ}; 180^{\circ}\right) \), pre ktoré \( \mathrm{tg}\,\alpha = \mathrm{cotg}\,\alpha \)?}
{\( 2 \)}{\( 1 \)}{\( 0 \)}{\( 4 \)}{}{}}
\LANGes{
\Question{
¿Cuántos valores del ángulo \( \alpha\in\left(0^{\circ}; 90^{\circ}\right)\cup\left(90^{\circ}; 180^{\circ}\right) \)  satisfacen la ecuación \( \mathrm{tg}\,\alpha = \mathrm{cotg}\,\alpha \)?}
{\( 2 \)}{\( 1 \)}{\( 0 \)}{\( 4 \)}{}{}}
\End

\Data\ID{1003076707}
\Updated{2020-07-15 03:07:12}
\Level{B}
\SubArea{Sine, cosine, tangent and cotangent}
\LANGen{
\Question{
How many values of the angle \( \alpha \in\left[0^{\circ}; 360^{\circ}\right] \) satisfy the equation 
\( \sin \alpha = \cos\alpha \)?}
{\( 2 \)}{\( 1 \)}{\( 0 \)}{\( 3 \)}{}{}}
\LANGpl{
\Question{
Jakie wartości kąta  \( \alpha \in\left\langle0^{\circ}; 360^{\circ}\right\rangle \) spełniają równanie
\( \sin \alpha = \cos\alpha \)?}
{\( 2 \)}{\( 1 \)}{\( 0 \)}{\( 3 \)}{}{}}
\LANGcs{
\Question{
Kolik existuje úhlů \( \alpha \in\left\langle0^{\circ}; 360^{\circ}\right\rangle \), pro které
\( \sin \alpha = \cos\alpha \)?}
{\( 2 \)}{\( 1 \)}{\( 0 \)}{\( 3 \)}{}{}}
\LANGsk{
\Question{
Koľko je takých uhlov \( \alpha \in\left\langle0^{\circ}; 360^{\circ}\right\rangle \), pre ktoré
\( \sin \alpha = \cos\alpha \)?}
{\( 2 \)}{\( 1 \)}{\( 0 \)}{\( 3 \)}{}{}}
\LANGes{
\Question{
¿Cuántos valores del ángulo \( \alpha \in\left[0^{\circ}; 360^{\circ}\right] \)  satisfacen la ecuación
\( \sin \alpha = \cos\alpha \)?}
{\( 2 \)}{\( 1 \)}{\( 0 \)}{\( 3 \)}{}{}}
\End

\Data\ID{1103076609}
\Updated{2020-07-15 03:07:12}
\Level{B}
\SubArea{Sine, cosine, tangent and cotangent}
\IMG{1103076609.pdf}
\LANGen{
\Question{
The graph in the picture is graph of  the function:}
{\( f(x)=-\sin\left(x + \frac{\pi}4 \right) \)}{\( f(x)=-\sin\left(x -\frac{3\pi}4 \right) \)}{\( f(x)=\cos\left(x + \frac{\pi}4\right) \)}{\( f(x)=\sin\left(x - \frac{\pi}4\right) \)}{}{}}
\LANGpl{
\Question{
Wskaż funkcję, której wykres jest przedstawiony na rysunku:}
{\( f(x)=-\sin\left(x + \frac{\pi}4 \right) \)}{\( f(x)=-\sin\left(x -\frac{3\pi}4 \right) \)}{\( f(x)=\cos\left(x + \frac{\pi}4\right) \)}{\( f(x)=\sin\left(x - \frac{\pi}4\right) \)}{}{}}
\LANGcs{
\Question{
Graf na obrázku patří funkci:}
{\( f(x)=-\sin\left(x + \frac{\pi}4 \right) \)}{\( f(x)=-\sin\left(x -\frac{3\pi}4 \right) \)}{\( f(x)=\cos\left(x + \frac{\pi}4\right) \)}{\( f(x)=\sin\left(x - \frac{\pi}4\right) \)}{}{}}
\LANGsk{
\Question{
Graf na obrázku patrí funkcii:}
{\( f(x)=-\sin\left(x + \frac{\pi}4 \right) \)}{\( f(x)=-\sin\left(x -\frac{3\pi}4 \right) \)}{\( f(x)=\cos\left(x + \frac{\pi}4\right) \)}{\( f(x)=\sin\left(x - \frac{\pi}4\right) \)}{}{}}
\LANGes{
\Question{
La gráfica mostrada en la imagen es la gráfica de la función:}
{\( f(x)=-\sin\left(x + \frac{\pi}4 \right) \)}{\( f(x)=-\sin\left(x -\frac{3\pi}4 \right) \)}{\( f(x)=\cos\left(x + \frac{\pi}4\right) \)}{\( f(x)=\sin\left(x - \frac{\pi}4\right) \)}{}{}}
\End

\Data\ID{1103076612}
\Updated{2020-07-15 03:07:12}
\Level{B}
\SubArea{Sine, cosine, tangent and cotangent}
\IMG{1103076612.pdf}
\LANGen{
\Question{
Determine the function  whose graph is given in the picture:}
{\( f(x)=-\cos 2x  \)}{\( f(x)=\cos(-2x) \)}{\( f(x)=-\sin 2x \)}{\( f(x)=\sin(-2x) \)}{}{}}
\LANGpl{
\Question{
Wskaż funkcję, której wykres jest przedstawiony na rysunku:}
{\( f(x)=-\cos 2x  \)}{\( f(x)=\cos(-2x) \)}{\( f(x)=-\sin 2x \)}{\( f(x)=\sin(-2x) \)}{}{}}
\LANGcs{
\Question{
Na obrázku je část grafu funkce:}
{\( f(x)=-\cos 2x  \)}{\( f(x)=\cos(-2x) \)}{\( f(x)=-\sin 2x \)}{\( f(x)=\sin(-2x) \)}{}{}}
\LANGsk{
\Question{
Na obrázku je časť grafu funkcie:}
{\( f(x)=-\cos 2x  \)}{\( f(x)=\cos(-2x) \)}{\( f(x)=-\sin 2x \)}{\( f(x)=\sin(-2x) \)}{}{}}
\LANGes{
\Question{
Determina la función cuya gráfica es la siguiente:}
{\( f(x)=-\cos 2x  \)}{\( f(x)=\cos(-2x) \)}{\( f(x)=-\sin 2x \)}{\( f(x)=\sin(-2x) \)}{}{}}
\End

\Data\ID{1103076613}
\Updated{2021-04-24 04:04:53}
\Level{B}
\SubArea{Sine, cosine, tangent and cotangent}
\LANGen{
\Question{
Which of the graphs is graph of   the function \( f(x)=2\cos(3x+\pi) \)?}
{}{}{}{}{}{}}
\LANGpl{
\Question{
Wskaż wykres funkcji \( f(x)=2\cos(3x+\pi) \)?}
{}{}{}{}{}{}}
\LANGcs{
\Question{
Který z grafů je grafem funkce \( f(x)=2\cos(3x+\pi) \)?}
{}{}{}{}{}{}}
\LANGsk{
\Question{
Ktorý z grafov je graf funkcie \( f(x)=2\cos(3x+\pi) \)?}
{}{}{}{}{}{}}
\LANGes{
\Question{
NA}
{}{}{}{}{}{}}

\IMGanswer{1103076613a.pdf}
\IMGanswer{1103076613b.pdf}
\IMGanswer{1103076613c.pdf}
\IMGanswer{1103076613d.pdf}\End

\Data\ID{1103076614}
\Updated{2020-07-15 03:07:10}
\Level{B}
\SubArea{Sine, cosine, tangent and cotangent}
\LANGen{
\Question{
Determine the graph of the function \( f(x)=0.5\sin(2x) + 1 \).}
{}{}{}{}{}{}}
\LANGpl{
\Question{
Wskaż wykres funkcji \( f(x)=0{,}5\sin(2x) + 1 \).}
{}{}{}{}{}{}}
\LANGcs{
\Question{
Který z grafů je grafem funkce \( f(x)=0{,}5\sin(2x) + 1 \)?}
{}{}{}{}{}{}}
\LANGsk{
\Question{
Ktorý z grafov je graf funkcie  \( f(x)=0{,}5\sin(2x) + 1 \)?}
{}{}{}{}{}{}}
\LANGes{
\Question{
NA}
{}{}{}{}{}{}}

\IMGanswer{1103076614a_0.pdf}
\IMGanswer{1103076614b_0.pdf}
\IMGanswer{1103076614c_0.pdf}
\IMGanswer{1103076614d_0.pdf}\End

\Data\ID{1103076615}
\Updated{2020-07-15 03:07:10}
\Level{B}
\SubArea{Sine, cosine, tangent and cotangent}
\LANGen{
\Question{
Choose the graph of the function \( f(x) = 3\cos\left(x + \frac{\pi}4 \right) \).}
{}{}{}{}{}{}}
\LANGpl{
\Question{
Wskaż wykres funkcji \( f(x) = 3\cos\left(x + \frac{\pi}4 \right) \).}
{}{}{}{}{}{}}
\LANGcs{
\Question{
Který z grafů je grafem funkce \( f(x) = 3\cos\left(x + \frac{\pi}4 \right) \)?}
{}{}{}{}{}{}}
\LANGsk{
\Question{
Vyberte graf funkcie \( f(x) = 3\cos\left(x + \frac{\pi}4 \right) \).}
{}{}{}{}{}{}}
\LANGes{
\Question{
NA}
{}{}{}{}{}{}}

\IMGanswer{1103076615a.pdf}
\IMGanswer{1103076615b.pdf}
\IMGanswer{1103076615c.pdf}
\IMGanswer{1103076615d.pdf}\End

\Data\ID{9000004801}
\Updated{2020-07-15 03:07:34}
\Level{B}
\SubArea{Sine, cosine, tangent and cotangent}
\LANGen{
\Question{
In the following list identify an even function.}
{\(y =\cos x\)}{\(y =\sin x\)}{\(y =\mathop{\mathrm{tg}}\nolimits x\)}{\(y =\mathop{\mathrm{cotg}}\nolimits x\)}{}{}}
\LANGpl{
\Question{
Która z niżej podanych funkcji jest funkcją parzystą?}
{\(y =\cos x\)}{\(y =\sin x\)}{\(y =\mathop{\mathrm{tg}}\nolimits x\)}{\(y =\mathop{\mathrm{cotg}}\nolimits x\)}{}{}}
\LANGcs{
\Question{
Která z uevedených funkcí je funkce sudá?}
{\(y =\cos x\)}{\(y =\sin x\)}{\(y =\mathop{\mathrm{tg}}\nolimits x\)}{\(y =\mathop{\mathrm{cotg}}\nolimits x\)}{}{}}
\LANGsk{
\Question{
Ktorá z uevedených funkcií je funkcia párna?}
{\(y =\cos x\)}{\(y =\sin x\)}{\(y =\mathop{\mathrm{tg}}\nolimits x\)}{\(y =\mathop{\mathrm{cotg}}\nolimits x\)}{}{}}
\LANGes{
\Question{
En la siguiente lista, identifica una función par.}
{\(y =\cos x\)}{\(y =\sin x\)}{\(y =\mathop{\mathrm{tg}}\nolimits x\)}{\(y =\mathop{\mathrm{cotg}}\nolimits x\)}{}{}}
\End

\Data\ID{9000004807}
\Updated{2020-07-15 03:07:34}
\Level{B}
\SubArea{Sine, cosine, tangent and cotangent}
\LANGen{
\Question{
In the following list identify a bounded function.}
{\(y =\sin x\)}{\(y =\mathop{\mathrm{tg}}\nolimits x\)}{\(y =\mathop{\mathrm{cotg}}\nolimits x\)}{\(y = -\mathop{\mathrm{tg}}\nolimits x\)}{}{}}
\LANGpl{
\Question{
Która z poniższych funkcji jest funkcją ograniczoną?}
{\(y =\sin x\)}{\(y =\mathop{\mathrm{tg}}\nolimits x\)}{\(y =\mathop{\mathrm{cotg}}\nolimits x\)}{\(y = -\mathop{\mathrm{tg}}\nolimits x\)}{}{}}
\LANGcs{
\Question{
Která z níže uvedených funkcí je ohraničená?}
{\(y =\sin x\)}{\(y =\mathop{\mathrm{tg}}\nolimits x\)}{\(y =\mathop{\mathrm{cotg}}\nolimits x\)}{\(y = -\mathop{\mathrm{tg}}\nolimits x\)}{}{}}
\LANGsk{
\Question{
Ktorá z nižšie uvedených funkcií je ohraničená?}
{\(y =\sin x\)}{\(y =\mathop{\mathrm{tg}}\nolimits x\)}{\(y =\mathop{\mathrm{cotg}}\nolimits x\)}{\(y = -\mathop{\mathrm{tg}}\nolimits x\)}{}{}}
\LANGes{
\Question{
En la siguiente lista, identifica una función que está acotada.}
{\(y =\sin x\)}{\(y =\mathop{\mathrm{tg}}\nolimits x\)}{\(y =\mathop{\mathrm{cotg}}\nolimits x\)}{\(y = -\mathop{\mathrm{tg}}\nolimits x\)}{}{}}
\End

\Data\ID{9000033801}
\Updated{2020-07-15 03:07:34}
\Level{B}
\SubArea{Sine, cosine, tangent and cotangent}
\LANGen{
\Question{
Which of the numbers in the list is a period (not necessarily the smallest period) of the
function \(m\colon y =\cos x\)?}
{\(4\pi \)}{\(\pi \)}{\(5\pi \)}{\(3\pi \)}{}{}}
\LANGpl{
\Question{
Która z podanych liczb jest okresem funkcji \(m\colon y =\cos x\)?}
{\(4\pi \)}{\(\pi \)}{\(5\pi \)}{\(3\pi \)}{}{}}
\LANGcs{
\Question{
Které z nabídnutých čísel lze považovat za periodu funkce
\(m\colon y =\cos x\)?}
{\(4\pi \)}{\(\pi \)}{\(5\pi \)}{\(3\pi \)}{}{}}
\LANGsk{
\Question{
Ktoré z ponúknutých čísel možno považovať za periódu funkcie
\(m\colon y =\cos x\)?}
{\(4\pi \)}{\(\pi \)}{\(5\pi \)}{\(3\pi \)}{}{}}
\LANGes{
\Question{
¿Cuál de los números en la lista es un período (no necesariamente el período más pequeño) de la
función \(m\colon y =\cos x\)?}
{\(4\pi \)}{\(\pi \)}{\(5\pi \)}{\(3\pi \)}{}{}}
\End

\Data\ID{9000033802}
\Updated{2020-07-15 03:07:34}
\Level{B}
\SubArea{Sine, cosine, tangent and cotangent}
\LANGen{
\Question{
Which of the numbers in the list is a period (not necessarily the smallest period) of the
function \(n\colon y =\mathop{\mathrm{tg}}\nolimits x\)?}
{\(3\pi \)}{\(\frac{\pi }{2}\)}{\(- \frac{\pi } {2}\)}{\(\frac{3\pi }
{2}\)}{}{}}
\LANGpl{
\Question{
Która z podanych liczb jest okresem funkcji \(n\colon y =\mathop{\mathrm{tg}}\nolimits x\)?}
{\(3\pi \)}{\(\frac{\pi }{2}\)}{\(- \frac{\pi } {2}\)}{\(\frac{3\pi }
{2}\)}{}{}}
\LANGcs{
\Question{
Které z nabídnutých čísel lze považovat za periodu funkce
\(n\colon y =\mathop{\mathrm{tg}}\nolimits x\)?}
{\(3\pi \)}{\(\frac{\pi }{2}\)}{\(- \frac{\pi } {2}\)}{\(\frac{3\pi }
{2}\)}{}{}}
\LANGsk{
\Question{
Ktoré z ponúknutých čísel možno považovať za periódu funkcie
\(n\colon y =\mathop{\mathrm{tg}}\nolimits x\)?}
{\(3\pi \)}{\(\frac{\pi }{2}\)}{\(- \frac{\pi } {2}\)}{\(\frac{3\pi }
{2}\)}{}{}}
\LANGes{
\Question{
¿Cuál de los números en la lista es un período (no necesariamente el período más pequeño) de la función \(n\colon y =\mathop{\mathrm{tg}}\nolimits x\)?}
{\(3\pi \)}{\(\frac{\pi }{2}\)}{\(- \frac{\pi } {2}\)}{\(\frac{3\pi }
{2}\)}{}{}}
\End

\Data\ID{9000033803}
\Updated{2021-04-24 04:04:32}
\Level{B}
\SubArea{Sine, cosine, tangent and cotangent}
\LANGen{
\Question{
In the following list identify a true statement for the function
\(f\colon y =\sin x\),
\(x\in \left [ -\frac{\pi }{2}; \frac{\pi } {2}\right ] \).}
{The function \(f\) 
is increasing.}{The function \(f\) 
is decreasing.}{The function \(f\) 
is neither increasing nor decreasing.}{}{}{}}
\LANGpl{
\Question{
Dana jest funkcja 
\(f\colon y =\sin x\),
\(x\in \left \langle -\frac{\pi }{2}; \frac{\pi } {2}\right \rangle \), wskaż zdanie prawdziwe.}
{Funkcja \(f\) 
jest rosnąca.}{Funkcja \(f\) 
jest malejąca.}{Funkcja  \(f\) 
nie jest ani rosnąca, ani malejąca.}{}{}{}}
\LANGcs{
\Question{
Je dána funkce \(f\colon y =\sin x\),
\(x\in \left \langle -\frac{\pi }{2}; \frac{\pi } {2}\right \rangle \).
Vyberte pravdivé tvrzení.}
{Funkce \(f\) 
je rostoucí.}{Funkce \(f\) 
je klesající.}{Funkce \(f\) 
není rostoucí, ani klesající.}{}{}{}}
\LANGsk{
\Question{
Je daná funkcia \(f\colon y =\sin x\),
\(x\in \left \langle -\frac{\pi }{2}; \frac{\pi } {2}\right \rangle \).
Vyberte pravdivé tvrdenie.}
{Funkcia \(f\) 
je rastúca.}{Funkcia \(f\) 
je klesajúca.}{Funkcia \(f\) 
nie je rastúca, ani klesajúca.}{}{}{}}
\LANGes{
\Question{
En la siguiente lista, identifica una proposición verdadera sobre la función
\(f\colon y =\sin x\),
\(x\in \left [ -\frac{\pi }{2}; \frac{\pi } {2}\right ] \).}
{La función \(f\) 
es creciente.}{La función \(f\) 
es decreciente.}{La función \(f\) 
no es ni creciente ni decreciente.}{}{}{}}
\End

\Data\ID{9000033804}
\Updated{2020-07-15 03:07:34}
\Level{B}
\SubArea{Sine, cosine, tangent and cotangent}
\LANGen{
\Question{
In the following list identify a true statement for the function
\(g\colon y =\sin x\),
\(x\in [ - 2\pi ;-\pi ] \).}
{The function \(g\) 
is neither increasing nor decreasing.}{The function \(g\) 
is increasing.}{The function \(g\) 
is decreasing.}{}{}{}}
\LANGpl{
\Question{
Dana jest funkcja
\(g\colon y =\sin x\),
\(x\in \langle - 2\pi ;-\pi \rangle \), wskaż zdanie prawdzie.}
{Funkcja  \(g\) 
nie jest ani rosnąca, ani malejąca.}{Funkcja  \(g\) 
jest rosnąca.}{Funkcja \(g\) 
malejąca.}{}{}{}}
\LANGcs{
\Question{
Je dána funkce \(g\colon y =\sin x\),
\(x\in \langle - 2\pi ;-\pi \rangle \).
Vyberte pravdivé tvrzení.}
{Funkce \(g\) 
není rostoucí, ani klesající.}{Funkce \(g\) 
je rostoucí.}{Funkce \(g\) 
je klesající.}{}{}{}}
\LANGsk{
\Question{
Je daná funkcia \(g\colon y =\sin x\),
\(x\in \langle - 2\pi ;-\pi \rangle \).
Vyberte pravdivé tvrdenie.}
{Funkcia \(g\) 
nie je rastúca, ani klesajúca.}{Funkcia \(g\) 
je rastúca.}{Funkcia \(g\) 
je klesajúca.}{}{}{}}
\LANGes{
\Question{
En la siguiente lista, identifica una proposición verdadera sobre la función
\(g\colon y =\sin x\),
\(x\in [ - 2\pi ;-\pi ] \).}
{La función \(g\) no es ni creciente ni decreciente.}{La función \(g\) 
es creciente.}{La función \(g\) 
es decreciente.}{}{}{}}
\End

\Data\ID{9000033805}
\Updated{2020-07-15 03:07:34}
\Level{B}
\SubArea{Sine, cosine, tangent and cotangent}
\LANGen{
\Question{
In the following list identify a true statement for the function
\(h\colon y =\mathop{\mathrm{cotg}}\nolimits x\),
\(x\in \left (-\frac{\pi }{2};0\right )\cup \left (0; \frac{\pi } {2}\right )\).}
{The function \(h\) 
is neither increasing nor decreasing.}{The function \(h\) 
is increasing.}{The function \(h\) 
is decreasing.}{}{}{}}
\LANGpl{
\Question{
Dana jest funkcja
\(h\colon y =\mathop{\mathrm{cotg}}\nolimits x\),
\(x\in \left (-\frac{\pi }{2};0\right )\cup \left (0;\frac{\pi } {2}\right )\), wskaż zdanie prawdziwe.}
{Funkcja \(h\) 
nie jest ani rosnąca, ani malejąca.}{Funkcja \(h\) 
jest rosnąca.}{Funkcja  \(h\) 
jest malejąca.}{}{}{}}
\LANGcs{
\Question{
Je dána funkce \(h\colon y =\mathop{\mathrm{cotg}}\nolimits x\),
\(x\in \left (-\frac{\pi }{2};0\right )\cup \left (0; \frac{\pi } {2}\right )\).
Vyberte pravdivé tvrzení.}
{Funkce \(h\) 
není rostoucí, ani klesající.}{Funkce \(h\) 
je rostoucí.}{Funkce \(h\) 
je klesající.}{}{}{}}
\LANGsk{
\Question{
Je daná funkcia \(h\colon y =\mathop{\mathrm{cotg}}\nolimits x\),
\(x\in \left (-\frac{\pi }{2};0\right )\cup \left (0; \frac{\pi } {2}\right )\).
Vyberte pravdivé tvrdenie.}
{Funkcia \(h\) 
nie je rastúca, ani klesajúca.}{Funkcia \(h\) 
je rastúca.}{Funkcia \(h\) 
je klesajúca.}{}{}{}}
\LANGes{
\Question{
En la siguiente lista, identifica una proposición verdadera sobre la función
\(h\colon y =\mathop{\mathrm{cotg}}\nolimits x\),
\(x\in \left (-\frac{\pi }{2};0\right )\cup \left (0; \frac{\pi } {2}\right )\).}
{La función \(h\) 
no es ni creciente ni decreciente.}{La función  \(h\) 
es creciente.}{La función \(h\) 
es decreciente.}{}{}{}}
\End

\Data\ID{9000033806}
\Updated{2020-07-15 03:07:34}
\Level{B}
\SubArea{Sine, cosine, tangent and cotangent}
\LANGen{
\Question{
In the following list identify a true statement for the function
\(i\colon y =\mathop{\mathrm{tg}}\nolimits x\),
\(x\in \left ( \frac{\pi }{2}; \frac{3\pi } 
{2}\right )\).}
{The function \(i\) 
is increasing.}{The function \(i\) 
is decreasing.}{The function \(i\) 
is neither increasing nor decreasing.}{}{}{}}
\LANGpl{
\Question{
Dana jest funkcja
\(i\colon y =\mathop{\mathrm{tg}}\nolimits x\),
\(x\in \left ( \frac{\pi }{2}; \frac{3\pi } 
{2}\right )\), wskaż zdanie prawdziwe.}
{Funkcja \(i\) 
jest rosnąca.}{Funkcja  \(i\) 
jest malejąca.}{Funkcja \(i\) 
nie jest ani rosnąca, ani malejąca.}{}{}{}}
\LANGcs{
\Question{
Je dána funkce \(i\colon y =\mathop{\mathrm{tg}}\nolimits x\),
\(x\in \left ( \frac{\pi }{2}; \frac{3\pi } 
{2}\right )\).
Vyberte pravdivé tvrzení.}
{Funkce \(i\) 
je rostoucí.}{Funkce \(i\) 
je klesající.}{Funkce \(i\) 
není rostoucí, ani klesající.}{}{}{}}
\LANGsk{
\Question{
Je daná funkcia \(i\colon y =\mathop{\mathrm{tg}}\nolimits x\),
\(x\in \left ( \frac{\pi }{2}; \frac{3\pi } 
{2}\right )\).
Vyberte pravdivé tvrdenie.}
{Funkcia \(i\) 
je rastúca.}{Funkcia \(i\) 
je klesajúca.}{Funkcia \(i\) 
nie je rastúca, ani klesajúca.}{}{}{}}
\LANGes{
\Question{
En la siguiente lista, identifica una proposición verdadera sobre la función 
\(i\colon y =\mathop{\mathrm{tg}}\nolimits x\),
\(x\in \left ( \frac{\pi }{2}; \frac{3\pi } 
{2}\right )\).}
{La función \(i\) 
es creciente.}{La función \(i\) 
es decreciente.}{La función \(i\) 
no es ni creciente ni decreciente.}{}{}{}}
\End

\Data\ID{9000033807}
\Updated{2021-04-24 04:04:03}
\Level{B}
\SubArea{Sine, cosine, tangent and cotangent}
\LANGen{
\Question{
In the following list identify a true statement for the function
\(f\colon y =\cos x\) on the
interval \(I = \left (-\frac{\pi }{2}; \frac{\pi } {2}\right )\).}
{The function \(f\) 
possess a unique maximum and no minimum on
\(I\).}{The function \(f\) does not have
a minimum or maximum on \(I\).}{The function \(f\) 
possess a unique maximum and unique minimum on
\(I\).}{The function \(f\) 
possess a unique minimum and no maximum on
\(I\).}{}{}}
\LANGpl{
\Question{
Dana jest funkcja 
\(f\colon y =\cos x\) na przedziale 
 \(I = \left (-\frac{\pi }{2}; \frac{\pi } {2}\right )\), wskaż zdanie prawdziwe.}
{Funkcja  \(f\) 
ma tylko jedno maksimum, nie ma minimum na przedziale
\(I\).}{Funkcja  \(f\) nie ma 
 minimum, ani maksimum na przedziale \(I\).}{Funkcja  \(f\) 
ma tylko jedno  maksimum i minimum na przedziale
\(I\).}{Funkcja  \(f\) 
ma tylko jedno minimum, nie ma maksimum na przedziale
\(I\).}{}{}}
\LANGcs{
\Question{
Pro extrémy funkce \(f\colon y =\cos x\) 
v intervalu \(\left (-\frac{\pi }{2}; \frac{\pi } {2}\right )\) 
platí:}
{V tomto intervalu existuje jediné maximum funkce
\(f\) a minimum
funkce \(f\) 
neexistuje.}{V tomto intervalu funkce \(f\) 
nemá žádný extrém.}{V tomto intervalu existuje jediné maximum a jediné minimum funkce
\(f\).}{V tomto intervalu existuje jediné minimum funkce
\(f\) a maximum
funkce \(f\) 
neexistuje.}{}{}}
\LANGsk{
\Question{
Pre extrémy funkcie \(f\colon y =\cos x\) 
v intervale \(\left (-\frac{\pi }{2}; \frac{\pi } {2}\right )\) 
platí:}
{V tomto intervale existuje jediné maximum funkcie
\(f\) a minimum
funkcie \(f\) 
neexistuje.}{V tomto intervale funkcia \(f\) 
nemá žiadny extrém.}{V tomto intervale existuje jediné maximum a jediné minimum funkcie
\(f\).}{V tomto intervale existuje jediné minimum funkcie
\(f\) a maximum
funkcie \(f\) 
neexistuje.}{}{}}
\LANGes{
\Question{
En la siguiente lista, identifica una proposición verdadera sobre la función 
\(f\colon y =\cos x\) en el intervalo \(I = \left (-\frac{\pi }{2}; \frac{\pi } {2}\right )\).}
{La función \(f\) 
posee un único máximo y ningún mínimo en
\(I\).}{La función  \(f\) no tiene
 mínimo o máximo en \(I\).}{La función  \(f\) 
posee un único máximo y un único mínimo en
\(I\).}{La función \(f\) 
posee un único mínimo y ningún máximo en
\(I\).}{}{}}
\End

\Data\ID{9000033808}
\Updated{2021-04-24 04:04:33}
\Level{B}
\SubArea{Sine, cosine, tangent and cotangent}
\LANGen{
\Question{
In the following list identify a true statement for the function
\(f\colon y =\sin x\) on the
interval \(I = \left (-\frac{\pi }{2}; \frac{\pi } {2}\right )\).}
{The function \(f\) does not have
a minimum or maximum on \(I\).}{The function \(f\) 
possess a unique maximum and unique minimum on
\(I\).}{The function \(f\) 
possess a unique maximum and no minimum on
\(I\).}{The function \(f\) 
possess a unique minimum and no maximum on
\(I\).}{}{}}
\LANGpl{
\Question{
Dana jest funkcja 
\(f\colon y =\sin x\) na przedziale 
 \(I = \left (-\frac{\pi }{2}; \frac{\pi } {2}\right )\), wskaż zdanie prawdziwe.}
{Funkcja  \(f\) nie ma 
 minimum, ani maksimum na przedziale \(I\).}{Funkcja \(f\) 
ma tylko jedno maksimum i minimum na przedziale
\(I\).}{Funkcja \(f\) 
ma tylko jedno  maksimum,  nie ma  minimum na przedziale
\(I\).}{Funkcja \(f\) 
ma tylko jedno minimum, nie ma  maksimum na przedziale
\(I\).}{}{}}
\LANGcs{
\Question{
Pro extrémy funkce \(f\colon y =\sin x\) 
v intervalu \(\left (-\frac{\pi }{2}; \frac{\pi } {2}\right )\) 
platí:}
{V tomto intervalu funkce \(f\) 
nemá žádný extrém.}{V tomto intervalu existuje jediné maximum a jediné minimum funkce
\(f\).}{V tomto intervalu existuje jediné maximum funkce
\(f\) a minimum
funkce \(f\) 
neexistuje.}{V tomto intervalu existuje jediné minimum funkce
\(f\) a maximum
funkce \(f\) 
neexistuje.}{}{}}
\LANGsk{
\Question{
Pre extrémy funkcie \(f\colon y =\sin x\) 
v intervale \(\left (-\frac{\pi }{2}; \frac{\pi } {2}\right )\) 
platí:}
{V tomto intervale funkcia \(f\) 
nemá žiadny extrém.}{V tomto intervale existuje jediné maximum a jediné minimum funkcie
\(f\).}{V tomto intervale existuje jediné maximum funkcie
\(f\) a minimum
funkcie \(f\) 
neexistuje.}{V tomto intervale existuje jediné minimum funkcie
\(f\) a maximum
funkcie \(f\) 
neexistuje.}{}{}}
\LANGes{
\Question{
En la siguiente lista, identifica una proposición verdadera sobre la función 
\(f\colon y =\sin x\) en el intervalo \(I = \left (-\frac{\pi }{2}; \frac{\pi } {2}\right )\).}
{La función  \(f\) no tiene
 mínimo o máximo en \(I\).}{La función  \(f\) 
posee un único máximo y un único mínimo en
\(I\).}{La función \(f\) 
posee un único máximo y ningún mínimo en
\(I\).}{La función  \(f\) 
posee un único mínimo y ningún máximo en
\(I\).}{}{}}
\End

\Data\ID{9000033809}
\Updated{2020-07-15 03:07:34}
\Level{B}
\SubArea{Sine, cosine, tangent and cotangent}
\LANGen{
\Question{
Find the parity (odd/even) of the function
\(k\colon y = -\mathop{\mathrm{tg}}\nolimits x\).}
{The function \(k\) 
is an odd function.}{The function \(k\) 
is an even function.}{The function \(k\) 
is neither an odd nor even function.}{}{}{}}
\LANGpl{
\Question{
Sprawdź parzystość  (tzn. parzysta/nieparzysta) funkcji
\(k\colon y = -\mathop{\mathrm{tg}}\nolimits x\).}
{Funkcja \(k\) 
jest nieparzysta.}{Funkcja  \(k\) 
jest parzysta.}{Funkcja \(k\) nie jest 
 parzysta, ani nieparzysta.}{}{}{}}
\LANGcs{
\Question{
Rozhodněte o paritě (tzn. o sudosti/lichosti) funkce
\(k\colon y = -\mathop{\mathrm{tg}}\nolimits x\).}
{Funkce \(k\) 
je lichá.}{Funkce \(k\) 
je sudá.}{Funkce \(k\) 
není ani sudá, ani lichá.}{}{}{}}
\LANGsk{
\Question{
Rozhodnite o párnosti alebo nepárnosti funkcie
\(k\colon y = -\mathop{\mathrm{tg}}\nolimits x\).}
{Funkcia \(k\) 
je nepárna.}{Funkcia \(k\) 
je párna.}{Funkcia \(k\) 
nie je ani párna, ani nepárna.}{}{}{}}
\LANGes{
\Question{
Determina la paridad (impar / par) de la función
\(k\colon y = -\mathop{\mathrm{tg}}\nolimits x\).}
{La función  \(k\) 
es una función impar.}{La función  \(k\) 
es una función par.}{La función   \(k\) 
no es una función impar ni par.}{}{}{}}
\End

\Data\ID{9000038901}
\Updated{2020-07-15 03:07:34}
\Level{B}
\SubArea{Sine, cosine, tangent and cotangent}
\LANGen{
\Question{
Consider the function \(f\colon y = A\cdot \sin (B\cdot x + C)\), with
real nonzero parameters \(A\),
\(B\) and
\(C\).
Which of the following operations makes the period of the function five times
smaller?}
{Increase \(B\) by
a factor \(5\).}{Increase \(A\) 
by a factor \(5\).}{Decrease \(A\) 
by a factor \(5\).}{Decrease \(B\) 
by a factor \(5\).}{Increase \(C\) 
by a factor \(5\).}{Decrease \(C\) 
by a factor \(5\).}}
\LANGpl{
\Question{
Dana jest funkcja \(f\colon y = A\cdot \sin (B\cdot x + C)\), gdzie
\(A\),
\(B\) i
\(C\) to rzeczywiste parametry niezerowe.
Które z poniższych działań zmniejszy okres funkcji pięciokrotnie?}
{Zwiększenie \(B\) pięć razy.}{Zwiększenie \(A\) pięć razy.}{Zmniejszenie \(A\) pięć razy.}{Zmniejszenie \(B\) pięć razy.}{Zwiększenie \(C\) pięć razy.}{Zmniejszenie \(C\) pięć razy.}}
\LANGcs{
\Question{
Je dána funkce \(f\colon y = A\cdot \sin (B\cdot x + C)\),
kde \(A\),
\(B\),
\(C\) jsou
reálné, nenulové parametry. Která z následujících změn parametru
pětkrát zmenší periodu funkce?}
{Pětkrát zvětšit \(B\).}{Pětkrát zvětšit \(A\).}{Pětkrát zmenšit \(A\).}{Pětkrát zmenšit \(B\).}{Pětkrát zvětšit \(C\).}{Pětkrát zmenšit \(C\).}}
\LANGsk{
\Question{
Je daná funkcia \(f\colon y = A\cdot \sin (B\cdot x + C)\),
kde \(A\),
\(B\),
\(C\) sú
reálne, nenulové parametre. Ktorá z nasledujúcich zmien parametra
päťkrát zmenší periódu funkcie?}
{Päťkrát zväčšiť \(B\).}{Päťkrát zväčšiť \(A\).}{Päťkrát zmenšiť \(A\).}{Päťkrát zmenšiť \(B\).}{Päťkrát zväčšiť \(C\).}{Päťkrát zmenšiť \(C\).}}
\LANGes{
\Question{
Considera la función \(f\colon y = A\cdot \sin (B\cdot x + C)\), con
parámetros reales distintos de cero \(A\),
\(B\) y
\(C\).
¿Cuál de las siguientes operaciones hace que la amplitud de la función sea cinco veces menor?}
{Aumentar  \(B\) en
un factor \(5\).}{Aumentar \(A\) 
en
un factor \(5\).}{Reducir  \(A\) 
en
un factor \(5\).}{Reducir \(B\) 
en
un factor \(5\).}{Aumentar \(C\) 
en
un factor \(5\).}{Reducir \(C\) 
en
un factor \(5\).}}
\End

\Data\ID{9000038902}
\Updated{2020-07-15 03:07:34}
\Level{B}
\SubArea{Sine, cosine, tangent and cotangent}
\LANGen{
\Question{
Consider the function \(f\colon y = A\cdot \sin (B\cdot x + C)\), with
real nonzero parameters \(A\),
\(B\) and
\(C\).
Which of the following operations makes the amplitude of the function five times
bigger?}
{Decrease \(A\) 
by a factor \(5\).}{Increase \(A\) by
a factor \(5\).}{Increase \(B\) 
by a factor \(5\).}{Decrease \(B\) 
by a factor \(5\).}{Increase \(C\) 
by a factor \(5\).}{Decrease \(C\) 
by a factor \(5\).}}
\LANGpl{
\Question{
Dana jest funkcja  \(f\colon y = A\cdot \sin (B\cdot x + C)\), gdzie
 \(A\),
\(B\) i
\(C\) to  rzeczywiste parametry niezerowe.
Które z poniższych działań zwiększy amplitudę  funkcji pięciokrotnie?}
{Zmniejszenie  \(A\) pięć razy.}{Zwiększenie \(A\) pięć razy.}{Zwiększenie \(B\) pięć razy.}{Zmniejszenie  \(B\) pięć razy.}{Zwiększenie \(C\) pięć razy.}{Zmniejszenie  \(C\) pięć razy.}}
\LANGcs{
\Question{
Je dána funkce \(f\colon y = A\cdot \sin (B\cdot x + C)\),
kde \(A\),
\(B\),
\(C\) jsou
reálné, nenulové parametry. Která z následujících změn parametru
pětkrát zmenší amplitudu funkce?}
{Pětkrát zmenšit \(A\).}{Pětkrát zvětšit \(A\).}{Pětkrát zvětšit \(B\).}{Pětkrát zmenšit \(B\).}{Pětkrát zvětšit \(C\).}{Pětkrát zmenšit \(C\).}}
\LANGsk{
\Question{
Je daná funkcia \(f\colon y = A\cdot \sin (B\cdot x + C)\),
kde \(A\),
\(B\),
\(C\) sú
reálne, nenulové parametre. Ktorá z nasledujúcich zmien parametra
päťkrát zmenší amplitúdu funkcie?}
{Päťkrát zmenšiť \(A\).}{Päťkrát zväčšiť \(A\).}{Päťkrát zväčšiť \(B\).}{Päťkrát zmenšiť \(B\).}{Päťkrát zväčšiť \(C\).}{Päťkrát zmenšiť \(C\).}}
\LANGes{
\Question{
Considera la función \(f\colon y = A\cdot \sin (B\cdot x + C)\),  con
parámetros reales distintos de cero \(A\),
\(B\) y
\(C\).
¿Cuál de las siguientes operaciones hace que la amplitud de la función sea cinco veces más grande?}
{Reducir \(A\) 
en
un factor \(5\).}{Aumentar \(A\) en
un factor \(5\).}{Aumentar \(B\) 
en
un factor \(5\).}{Reducir \(B\) 
en
un factor \(5\).}{Aumentar \(C\) 
en
un factor \(5\).}{Reducir \(C\) 
en
un factor \(5\).}}
\End

\Data\ID{9000038903}
\Updated{2020-07-15 03:07:34}
\Level{B}
\SubArea{Sine, cosine, tangent and cotangent}
\LANGen{
\Question{
Find the minimal value of the parameter
\(D\) which ensures
that the function \(f\colon y = D + 3\cdot \sin x\) 
is nonnegative.}
{\(D = 3\)}{\(D = -3\)}{\(D = 6\)}{\(D = -6\)}{\(D = 1\)}{No solution.}}
\LANGpl{
\Question{
Wyznacz wartość minimalną parametru 
\(D\) tak, aby funkcja
 \(f\colon y = D + 3\cdot \sin x\) 
była funkcją nieujemną.}
{\(D = 3\)}{\(D = -3\)}{\(D = 6\)}{\(D = -6\)}{\(D = 1\)}{Brak rozwiązania.}}
\LANGcs{
\Question{
Jakou minimální hodnotu musí mít parametr
\(D\), aby
funkce \(f\colon y = D + 3\cdot \sin x\) 
nabývala pouze nezáporných hodnot?}
{\(D = 3\)}{\(D = -3\)}{\(D = 6\)}{\(D = -6\)}{\(D = 1\)}{Tento parametr nemůže ovlivnit znaménko funkce.}}
\LANGsk{
\Question{
Akú minimálnu hodnotu musí mať parameter
\(D\), aby
funkcia \(f\colon y = D + 3\cdot \sin x\) 
nadobúdala iba nezáporné hodnoty?}
{\(D = 3\)}{\(D = -3\)}{\(D = 6\)}{\(D = -6\)}{\(D = 1\)}{Tento parameter nemôže ovplyvniť znamienko funkcie.}}
\LANGes{
\Question{
Determina el valor mínimo del parámetro
\(D\)  suponiendo que 
 la función \(f\colon y = D + 3\cdot \sin x\) 
no es negativa.}
{\(D = 3\)}{\(D = -3\)}{\(D = 6\)}{\(D = -6\)}{\(D = 1\)}{No hay solución.}}
\End

\Data\ID{9000038904}
\Updated{2021-04-24 04:04:49}
\Level{B}
\SubArea{Sine, cosine, tangent and cotangent}
\LANGen{
\Question{
Consider function \(f\colon y =\sin x\) with domain
restricted to \(\mathop{\mathrm{Dom}}(f) =\mathbb{R}_{ 0}^{+}\). Identify the
function which is defined on \([ - 5;+\infty )\).}
{\(f(x + 5)\)}{\(f(x - 5)\)}{\(f(x) + 5\)}{\(f(x) - 5\)}{\(5\cdot f(x)\)}{}}
\LANGpl{
\Question{
Dana jest funkcja \(f\colon y =\sin x\), dziedziną funkcji jest 
 \(\mathop{\mathrm{D}}(f) =\mathbb{R}_{ 0}^{+}\). Wyznacz funkcję 
określoną na przedziale \(\langle - 5;+\infty )\).}
{\(f(x + 5)\)}{\(f(x - 5)\)}{\(f(x) + 5\)}{\(f(x) - 5\)}{\(5\cdot f(x)\)}{}}
\LANGcs{
\Question{
Je dána funkce \(f\colon y =\sin x\) s
definičním oborem \(D(f) =\mathbb{R}_{ 0}^{+}\).
Určete, která z následujících funkcí má definiční obor
\(\langle - 5;+\infty )\).}
{\(f(x + 5)\)}{\(f(x - 5)\)}{\(f(x) + 5\)}{\(f(x) - 5\)}{\(5\cdot f(x)\)}{}}
\LANGsk{
\Question{
Je daná funkcia \(f\colon y =\sin x\) s
definičným oborom \(D(f) =\mathbb{R}_{ 0}^{+}\).
Určte, ktorá z nasledujúcich funkcií má definičný obor
\(\langle - 5;+\infty )\).}
{\(f(x + 5)\)}{\(f(x - 5)\)}{\(f(x) + 5\)}{\(f(x) - 5\)}{\(5\cdot f(x)\)}{}}
\LANGes{
\Question{
Considera la función \(f\colon y =\sin x\) con el dominio \(\mathop{\mathrm{Dom}}(f) =\mathbb{R}_{ 0}^{+}\). Identifica cuál de las siguientes funciones tiene como dominio 
 \([ - 5;+\infty )\).}
{\(f(x + 5)\)}{\(f(x - 5)\)}{\(f(x) + 5\)}{\(f(x) - 5\)}{\(5\cdot f(x)\)}{}}
\End

\Data\ID{9000038905}
\Updated{2021-04-24 04:04:35}
\Level{B}
\SubArea{Sine, cosine, tangent and cotangent}
\LANGen{
\Question{
Identify the transformation which transforms the graph of the function
\(g\colon y =\sin 3x\) to the graph
of the function \(f\colon y =\sin (3x + 5)\).}
{Shift graph of \(g\) 
by \(\frac{5}
{3}\) 
units to the left.}{Shift graph of \(g\) 
by \(5\) 
units to the right.}{Shift graph of \(g\) 
by \(5\) 
units to the left.}{Shift graph of \(g\) 
by \(3\) 
units to the right.}{Shift graph of \(g\) 
by \(3\) 
units to the left.}{Shift graph of \(g\) 
by \(\frac{5}
{3}\) 
units to the right.}}
\LANGpl{
\Question{
Jak uzyskać wykres funkcji  \(f\colon y =\sin (3x + 5)\) z wykresu funkcji
\(g\colon y =\sin 3x\)?}
{Należy przesunąć wykres \(g\) 
o \(\frac{5}
{3}\) 
jednostki w lewo.}{Należy przesunąć wykres\(g\) 
o \(5\) 
jednostek w prawo.}{Należy przesunąć wykres \(g\) 
o \(5\) 
jednostek w lewo.}{Należy przesunąć wykres \(g\) 
o \(3\) 
jednostki w prawo.}{Należy przesunąć wykres \(g\) 
o \(3\) 
jednostki w lewo.}{Należy przesunąć wykres \(g\) 
o \(\frac{5}
{3}\) 
jednostki w prawo.}}
\LANGcs{
\Question{
Jak získáme graf funkce \(f\colon y =\sin (3x + 5)\) 
z grafu funkce \(g\colon y =\sin 3x\)?}
{Graf funkce \(g\) 
posuneme o \(\frac{5}
{3}\) ve směru
záporné poloosy \(x\).}{Graf funkce \(g\) posuneme o 5
ve směru kladné poloosy \(x\).}{Graf funkce \(g\) posuneme o 5 ve
směru záporné poloosy \(x\).}{Graf funkce \(g\) posuneme o 3
ve směru kladné poloosy \(x\).}{Graf funkce \(g\) posuneme o 3 ve
směru záporné poloosy \(x\).}{Graf funkce \(g\) 
posuneme o \(\frac{5}
{3}\) ve směru
kladné poloosy \(x\).}}
\LANGsk{
\Question{
Ako získame graf funkcie \(f\colon y =\sin (3x + 5)\) 
z grafu funkcie \(g\colon y =\sin 3x\)?}
{Graf funkcie \(g\) 
posunieme o \(\frac{5}
{3}\) v smere
zápornej polosi \(x\).}{Graf funkcie \(g\) posunieme o 5
v smere kladnej polosi \(x\).}{Graf funkcie \(g\) posunieme o 5 v
smere zápornej polosi \(x\).}{Graf funkcie \(g\) posunieme o 3
v smere kladnej polosi \(x\).}{Graf funkcie \(g\) posunieme o 3 v
smere zápornej polosi \(x\).}{Graf funkcie \(g\) 
posunieme o \(\frac{5}
{3}\) v smere
kladnej polosi \(x\).}}
\LANGes{
\Question{
¿Cómo obtenemos la gráfica de la función 
\(f\colon y =\sin (3x + 5)\)
partiendo de la gráfica de la función \(g\colon y =\sin 3x\)?}
{Desplazando la gráfica de  \(g\) 
hacia izquierda  \(\frac{5}
{3}\) 
unidades.}{Desplazando la gráfica de \(g\) 
hacia derecha \(5\) 
unidades.}{Desplazando la gráfica de \(g\) 
hacia izquierda  \(5\) 
unidades.}{Desplazando la gráfica de \(g\) 
hacia derecha \(3\) 
unidades.}{Desplazando la gráfica de \(g\) 
hacia izquierda  \(3\) 
unidades.}{Desplazando la gráfica de \(g\) 
hacia derecha  \(\frac{5}
{3}\) 
unidades.}}
\End

\Data\ID{9000038906}
\Updated{2021-04-24 04:04:07}
\Level{B}
\SubArea{Sine, cosine, tangent and cotangent}
\LANGen{
\Question{
Consider the function \(f\colon y =\mathop{\mathrm{tg}}\nolimits x\).
In the following list identify the nonnegative function.}
{None of the given functions is nonnegative.}{\(A\cdot f(x)\),
where \(A\in (-\infty ;0)\)}{\(A\cdot f(x)\),
where \(A\in (0;+\infty )\)}{\(f(B\cdot x)\),
where \(B\in (0;+\infty )\)}{\(f(x + C)\),
where \(C\in (-\infty ;0)\)}{}}
\LANGpl{
\Question{
Dana jest funkcja \(f\colon y =\mathop{\mathrm{tg}}\nolimits x\).
Wskaż funkcję nieujemną.}
{Brak funkcji nieujemnych.}{\(A\cdot f(x)\),
gdzie \(A\in (-\infty ;0)\)}{\(A\cdot f(x)\),
gdzie \(A\in (0;+\infty )\)}{\(f(B\cdot x)\),
gdzie \(B\in (0;+\infty )\)}{\(f(x + C)\),
gdzie \(C\in (-\infty ;0)\)}{}}
\LANGcs{
\Question{
Je dána funkce \(f\colon y =\mathop{\mathrm{tg}}\nolimits x\).
Určete, která z následujících funkcí bude mít pouze nezáporné hodnoty.}
{Žádná z daných uvedených funkcí nebude mít nezáporné hodnoty.}{\(A\cdot f(x)\), kde
\(A\in (-\infty ;0)\)}{\(A\cdot f(x)\), kde
\(A\in (0;+\infty )\)}{\(f(B\cdot x)\), kde
\(B\in (0;+\infty )\)}{\(f(x + C)\), kde
\(C\in (-\infty ;0)\)}{}}
\LANGsk{
\Question{
Je daná funkcia \(f\colon y =\mathop{\mathrm{tg}}\nolimits x\).
Určte, ktorá z daných funkcií bude mať iba nezáporné hodnoty.}
{Žiadna z vyššie uvedených funkcií nebude mať nezáporné hodnoty.}{\(A\cdot f(x)\), kde
\(A\in (-\infty ;0)\)}{\(A\cdot f(x)\), kde
\(A\in (0;+\infty )\)}{\(f(B\cdot x)\), kde
\(B\in (0;+\infty )\)}{\(f(x + C)\), kde
\(C\in (-\infty ;0)\)}{}}
\LANGes{
\Question{
Considera la función  \(f\colon y =\mathop{\mathrm{tg}}\nolimits x\).
En la siguiente lista, identifica la función no negativa.}
{Ninguna de las funciones dadas  es no negativa.}{\(A\cdot f(x)\),
donde \(A\in (-\infty ;0)\)}{\(A\cdot f(x)\),
donde \(A\in (0;+\infty )\)}{\(f(B\cdot x)\),
donde \(B\in (0;+\infty )\)}{\(f(x + C)\),
donde \(C\in (-\infty ;0)\)}{}}
\End

\Data\ID{9000038907}
\Updated{2021-04-24 04:04:42}
\Level{B}
\SubArea{Sine, cosine, tangent and cotangent}
\LANGen{
\Question{
Consider the function \(f\colon y =\mathop{\mathrm{cotg}}\nolimits x\) with
domain restricted to the interval \(\mathop{\mathrm{Dom}}(f) = (0;\pi )\).
In the following list identify the function with domain
\(\left (0; \frac{\pi } {3}\right )\).}
{\(f(3\cdot x)\)}{\(f(x - 3)\)}{\(f(x + 3)\)}{\(f\left (\frac{x}
{3} \right )\)}{\(3\cdot f(x)\)}{}}
\LANGpl{
\Question{
Dziedziną funkcji  \(f\colon y =\mathop{\mathrm{cotg}}\nolimits x\) jest przedział 
\(\mathop{\mathrm{D}}(f) = (0;\pi )\).
Wskaż funkcję, której dziedziną jest przedział
\(\left (0; \frac{\pi } {3}\right )\).}
{\(f(3\cdot x)\)}{\(f(x - 3)\)}{\(f(x + 3)\)}{\(f\left (\frac{x}
{3} \right )\)}{\(3\cdot f(x)\)}{}}
\LANGcs{
\Question{
Je dána funkce \(f\colon y =\mathop{\mathrm{cotg}}\nolimits x\) s
definičním oborem \(D(f) = (0;\pi )\).
Určete, která z následujících funkcí má definiční obor
\(\left (0; \frac{\pi } {3}\right )\).}
{\(f(3\cdot x)\)}{\(f(x - 3)\)}{\(f(x + 3)\)}{\(f\left (\frac{x}
{3} \right )\)}{\(3\cdot f(x)\)}{}}
\LANGsk{
\Question{
Je daná funkcia \(f\colon y =\mathop{\mathrm{cotg}}\nolimits x\) s
definičným oborom \(D(f) = (0;\pi )\).
Určte, ktorá z nasledujúcich funkcií má definičný obor
\(\left (0; \frac{\pi } {3}\right )\).}
{\(f(3\cdot x)\)}{\(f(x - 3)\)}{\(f(x + 3)\)}{\(f\left (\frac{x}
{3} \right )\)}{\(3\cdot f(x)\)}{}}
\LANGes{
\Question{
Considera la función  \(f\colon y =\mathop{\mathrm{cotg}}\nolimits x\) con dominio \(\mathop{\mathrm{Dom}}(f) = (0;\pi )\).
En la siguiente lista, identifica  la función con el dominio
\(\left (0; \frac{\pi } {3}\right )\).}
{\(f(3\cdot x)\)}{\(f(x - 3)\)}{\(f(x + 3)\)}{\(f\left (\frac{x}
{3} \right )\)}{\(3\cdot f(x)\)}{}}
\End

\Data\ID{9000038908}
\Updated{2021-04-24 04:04:10}
\Level{B}
\SubArea{Sine, cosine, tangent and cotangent}
\LANGen{
\Question{
Consider the function \(f\colon y =\mathop{\mathrm{tg}}\nolimits x\) with
domain restricted to the interval \(\mathop{\mathrm{Dom}}(f) = \left ( \frac{\pi }{2}; \frac{3\pi } 
{2}\right )\).
In the following list identify the function with domain
\((0;\pi )\).}
{\(f\left (x + \frac{\pi } {2}\right )\)}{\(\left ( \frac{\pi }{2}\right )\cdot f(x)\)}{\(f\left (x - \frac{\pi } {2}\right )\)}{\(f(x) + \frac{\pi } {2}\)}{\(f(x) - \frac{\pi } {2}\)}{}}
\LANGpl{
\Question{
Dziedziną funkcji \(f\colon y =\mathop{\mathrm{tg}}\nolimits x\) jest przedział  \(\mathop{\mathrm{D}}(f) = \left ( \frac{\pi }{2}; \frac{3\pi } 
{2}\right )\).
Wskaż funkcję, której dziedziną jest przedział
\((0;\pi )\).}
{\(f\left (x + \frac{\pi } {2}\right )\)}{\(\left ( \frac{\pi }{2}\right )\cdot f(x)\)}{\(f\left (x - \frac{\pi } {2}\right )\)}{\(f(x) + \frac{\pi } {2}\)}{\(f(x) - \frac{\pi } {2}\)}{}}
\LANGcs{
\Question{
Je dána funkce \(f\colon y =\mathop{\mathrm{tg}}\nolimits x\) s
definičním oborem \(D(f) = \left ( \frac{\pi }{2}; \frac{3\pi } 
{2}\right )\).
Určete, která z následujících funkcí má definiční obor
\((0;\pi )\).}
{\(f\left (x + \frac{\pi } {2}\right )\)}{\(\left ( \frac{\pi }{2}\right )\cdot f(x)\)}{\(f\left (x - \frac{\pi } {2}\right )\)}{\(f(x) + \frac{\pi } {2}\)}{\(f(x) - \frac{\pi } {2}\)}{}}
\LANGsk{
\Question{
Je daná funkcia \(f\colon y =\mathop{\mathrm{tg}}\nolimits x\) s
definičným oborom \(D(f) = \left ( \frac{\pi }{2}; \frac{3\pi } 
{2}\right )\).
Určte, ktorá z nasledujúcich funkcií má definičný obor
\((0;\pi )\).}
{\(f\left (x + \frac{\pi } {2}\right )\)}{\(\left ( \frac{\pi }{2}\right )\cdot f(x)\)}{\(f\left (x - \frac{\pi } {2}\right )\)}{\(f(x) + \frac{\pi } {2}\)}{\(f(x) - \frac{\pi } {2}\)}{}}
\LANGes{
\Question{
Considera la función \(f\colon y =\mathop{\mathrm{tg}}\nolimits x\) con dominio \(\mathop{\mathrm{Dom}}(f) = \left ( \frac{\pi }{2}; \frac{3\pi } 
{2}\right )\).
En la siguiente lista, identifica la función el dominio
\((0;\pi )\).}
{\(f\left (x + \frac{\pi } {2}\right )\)}{\(\left ( \frac{\pi }{2}\right )\cdot f(x)\)}{\(f\left (x - \frac{\pi } {2}\right )\)}{\(f(x) + \frac{\pi } {2}\)}{\(f(x) - \frac{\pi } {2}\)}{}}
\End

\Data\ID{9000038909}
\Updated{2020-07-15 03:07:34}
\Level{B}
\SubArea{Sine, cosine, tangent and cotangent}
\LANGen{
\Question{
Consider the function \(f\colon y =\sin \left (\frac{x}
{2} + \frac{\pi } 
{2}\right )\).
In the following list identify the function which has the same graph as the graph of the
function \(f\).}
{\(g\colon y =\cos \frac{x} 
{2} \)}{\(k\colon y =\cos \left (\frac{x}
{2} + \frac{\pi } 
{2}\right )\)}{\(b\colon y =\cos \left (\frac{x}
{2} - \frac{\pi } 
{2}\right )\)}{\(h\colon y =\cos \left (\frac{x}
{2} -\pi \right )\)}{\(m\colon y =\cos 2x\)}{}}
\LANGpl{
\Question{
Dana jest funkcja  \(f\colon y =\sin \left (\frac{x}
{2} + \frac{\pi } 
{2}\right )\).
Wskaż funkcję, która ma taki sam wykres jak funkcja \(f\).}
{\(g\colon y =\cos \frac{x} 
{2} \)}{\(k\colon y =\cos \left (\frac{x}
{2} + \frac{\pi } 
{2}\right )\)}{\(b\colon y =\cos \left (\frac{x}
{2} - \frac{\pi } 
{2}\right )\)}{\(h\colon y =\cos \left (\frac{x}
{2} -\pi \right )\)}{\(m\colon y =\cos 2x\)}{}}
\LANGcs{
\Question{
Určete, která z následujících funkcí má graf totožný s grafem funkce
\(f\colon y =\sin \left (\frac{x}
{2} + \frac{\pi } 
{2}\right )\).}
{\(g\colon y =\cos \frac{x} 
{2} \)}{\(k\colon y =\cos \left (\frac{x}
{2} + \frac{\pi } 
{2}\right )\)}{\(b\colon y =\cos \left (\frac{x}
{2} - \frac{\pi } 
{2}\right )\)}{\(h\colon y =\cos \left (\frac{x}
{2} -\pi \right )\)}{\(m\colon y =\cos 2x\)}{}}
\LANGsk{
\Question{
Určte, ktorá z nasledujúcich funkcií má graf totožný s grafom funkcie
\(f\colon y =\sin \left (\frac{x}
{2} + \frac{\pi } 
{2}\right )\).}
{\(g\colon y =\cos \frac{x} 
{2} \)}{\(k\colon y =\cos \left (\frac{x}
{2} + \frac{\pi } 
{2}\right )\)}{\(b\colon y =\cos \left (\frac{x}
{2} - \frac{\pi } 
{2}\right )\)}{\(h\colon y =\cos \left (\frac{x}
{2} -\pi \right )\)}{\(m\colon y =\cos 2x\)}{}}
\LANGes{
\Question{
Considera la función \(f\colon y =\sin \left (\frac{x}
{2} + \frac{\pi } 
{2}\right )\).
En la siguiente lista, identifica la función que tiene la misma gráfica que la
función \(f\).}
{\(g\colon y =\cos \frac{x} 
{2} \)}{\(k\colon y =\cos \left (\frac{x}
{2} + \frac{\pi } 
{2}\right )\)}{\(b\colon y =\cos \left (\frac{x}
{2} - \frac{\pi } 
{2}\right )\)}{\(h\colon y =\cos \left (\frac{x}
{2} -\pi \right )\)}{\(m\colon y =\cos 2x\)}{}}
\End

\Data\ID{9000038910}
\Updated{2020-07-15 03:07:34}
\Level{B}
\SubArea{Sine, cosine, tangent and cotangent}
\LANGen{
\Question{
Consider the function \(f\colon y =\mathop{\mathrm{cotg}}\nolimits x\).
In the following list identify the function which has the same graph as the graph of the
function \(f\).}
{\(k\colon y = -\mathop{\mathrm{tg}}\nolimits \left (x + \frac{\pi } {2}\right )\)}{\(g\colon y = -\mathop{\mathrm{tg}}\nolimits x\)}{\(b\colon y =\mathop{\mathrm{tg}}\nolimits \left (x + \frac{\pi } {2}\right )\)}{\(h\colon y =\mathop{\mathrm{tg}}\nolimits \left (x - \frac{\pi } {2}\right )\)}{\(m\colon y = -\mathop{\mathrm{tg}}\nolimits x - \frac{\pi } {2}\)}{}}
\LANGpl{
\Question{
Dana jest funkcja  \(f\colon y =\mathop{\mathrm{cotg}}\nolimits x\).
Wskaż funkcję, która ma taki sam wykres jak funkcja \(f\).}
{\(k\colon y = -\mathop{\mathrm{tg}}\nolimits \left (x + \frac{\pi } {2}\right )\)}{\(g\colon y = -\mathop{\mathrm{tg}}\nolimits x\)}{\(b\colon y =\mathop{\mathrm{tg}}\nolimits \left (x + \frac{\pi } {2}\right )\)}{\(h\colon y =\mathop{\mathrm{tg}}\nolimits \left (x - \frac{\pi } {2}\right )\)}{\(m\colon y = -\mathop{\mathrm{tg}}\nolimits x - \frac{\pi } {2}\)}{}}
\LANGcs{
\Question{
Určete, která z následujících funkcí má graf totožný s grafem funkce
\(f\colon y =\mathop{\mathrm{cotg}}\nolimits x\).}
{\(k\colon y = -\mathop{\mathrm{tg}}\nolimits \left (x + \frac{\pi } {2}\right )\)}{\(g\colon y = -\mathop{\mathrm{tg}}\nolimits x\)}{\(b\colon y =\mathop{\mathrm{tg}}\nolimits \left (x + \frac{\pi } {2}\right )\)}{\(h\colon y =\mathop{\mathrm{tg}}\nolimits \left (x - \frac{\pi } {2}\right )\)}{\(m\colon y = -\mathop{\mathrm{tg}}\nolimits x - \frac{\pi } {2}\)}{}}
\LANGsk{
\Question{
Určte, ktorá z nasledujúcich funkcií má graf totožný s grafom funkcie
\(f\colon y =\mathop{\mathrm{cotg}}\nolimits x\).}
{\(k\colon y = -\mathop{\mathrm{tg}}\nolimits \left (x + \frac{\pi } {2}\right )\)}{\(g\colon y = -\mathop{\mathrm{tg}}\nolimits x\)}{\(b\colon y =\mathop{\mathrm{tg}}\nolimits \left (x + \frac{\pi } {2}\right )\)}{\(h\colon y =\mathop{\mathrm{tg}}\nolimits \left (x - \frac{\pi } {2}\right )\)}{\(m\colon y = -\mathop{\mathrm{tg}}\nolimits x - \frac{\pi } {2}\)}{}}
\LANGes{
\Question{
Considera la función  \(f\colon y =\mathop{\mathrm{cotg}}\nolimits x\).
En la siguiente lista, identifica la función que tiene la misma gráfica que la función \(f\).}
{\(k\colon y = -\mathop{\mathrm{tg}}\nolimits \left (x + \frac{\pi } {2}\right )\)}{\(g\colon y = -\mathop{\mathrm{tg}}\nolimits x\)}{\(b\colon y =\mathop{\mathrm{tg}}\nolimits \left (x + \frac{\pi } {2}\right )\)}{\(h\colon y =\mathop{\mathrm{tg}}\nolimits \left (x - \frac{\pi } {2}\right )\)}{\(m\colon y = -\mathop{\mathrm{tg}}\nolimits x - \frac{\pi } {2}\)}{}}
\End

\Data\ID{1003076702}
\Updated{2020-07-15 03:07:12}
\Level{C}
\SubArea{Sine, cosine, tangent and cotangent}
\LANGen{
\Question{
The domain of the expression \( \frac{\cos x}{1-\sin x} \)  is the set:}
{\( \left\{x\in\mathbb{R}\colon x\neq\frac{\pi}2 + 2k\pi\text{, }  k\in\mathbb{Z} \right\} \)}{\( \mathbb{R} \)}{\( \left\{x\in\mathbb{R}\colon x\neq\frac{3\pi}2 + 2k\pi\text{, }  k\in\mathbb{Z} \right\} \)}{\( \left\{x\in\mathbb{R}\colon  x =k\pi\text{, } k\in\mathbb{Z} \right\} \)}{}{}}
\LANGpl{
\Question{
Dziedziną wyrażenia \( \frac{\cos x}{1-\sin x} \)  jest zbiór:}
{\( \left\{x\in\mathbb{R}\colon x\neq\frac{\pi}2 + 2k\pi\text{, }  k\in\mathbb{Z} \right\} \)}{\( \mathbb{R} \)}{\( \left\{x\in\mathbb{R}\colon x\neq\frac{3\pi}2 + 2k\pi\text{, }  k\in\mathbb{Z} \right\} \)}{\( \left\{x\in\mathbb{R}\colon  x =k\pi\text{, } k\in\mathbb{Z} \right\} \)}{}{}}
\LANGcs{
\Question{
Definičním oborem výrazu  \( \frac{\cos x}{1-\sin x} \)  je množina:}
{\( \left\{x\in\mathbb{R}\colon x\neq\frac{\pi}2 + 2k\pi\text{, }  k\in\mathbb{Z} \right\} \)}{\( \mathbb{R} \)}{\( \left\{x\in\mathbb{R}\colon x\neq\frac{3\pi}2 + 2k\pi\text{, }  k\in\mathbb{Z} \right\} \)}{\( \left\{x\in\mathbb{R}\colon  x =k\pi\text{, } k\in\mathbb{Z} \right\} \)}{}{}}
\LANGsk{
\Question{
Definičný obor výrazu    \( \frac{\cos x}{1-\sin x} \)  je množina:}
{\( \left\{x\in\mathbb{R}\colon x\neq\frac{\pi}2 + 2k\pi\text{, }  k\in\mathbb{Z} \right\} \)}{\( \mathbb{R} \)}{\( \left\{x\in\mathbb{R}\colon x\neq\frac{3\pi}2 + 2k\pi\text{, }  k\in\mathbb{Z} \right\} \)}{\( \left\{x\in\mathbb{R}\colon  x =k\pi\text{, } k\in\mathbb{Z} \right\} \)}{}{}}
\LANGes{
\Question{
El dominio de la expresión \( \frac{\cos x}{1-\sin x} \)  es el conjunto:}
{\( \left\{x\in\mathbb{R}\colon x\neq\frac{\pi}2 + 2k\pi\text{, }  k\in\mathbb{Z} \right\} \)}{\( \mathbb{R} \)}{\( \left\{x\in\mathbb{R}\colon x\neq\frac{3\pi}2 + 2k\pi\text{, }  k\in\mathbb{Z} \right\} \)}{\( \left\{x\in\mathbb{R}\colon  x =k\pi\text{, } k\in\mathbb{Z} \right\} \)}{}{}}
\End

\Data\ID{1003076703}
\Updated{2020-07-15 03:07:12}
\Level{C}
\SubArea{Sine, cosine, tangent and cotangent}
\LANGen{
\Question{
The expression \( \frac{\sin 2x}{\cos^2x } \) is equal to:}
{\( 2\,\mathrm{tg}\,x \)}{\( \frac{\sin x}{1-\sin x} \)}{\( \mathrm{tg}^2 x \)}{\( \mathrm{tg}\,2x \)}{}{}}
\LANGpl{
\Question{
Wyrażenie \( \frac{\sin 2x}{\cos^2x } \) jest równe:}
{\( 2\,\mathrm{tg}\,x \)}{\( \frac{\sin x}{1-\sin x} \)}{\( \mathrm{tg}^2 x \)}{\( \mathrm{tg}\,2x \)}{}{}}
\LANGcs{
\Question{
Výraz \( \frac{\sin 2x}{\cos^2x } \) je roven:}
{\( 2\,\mathrm{tg}\,x \)}{\( \frac{\sin x}{1-\sin x} \)}{\( \mathrm{tg}^2 x \)}{\( \mathrm{tg}\,2x \)}{}{}}
\LANGsk{
\Question{
Po úprave výrazu  \( \frac{\sin 2x}{\cos^2x } \)  dostaneme:}
{\( 2\,\mathrm{tg}\,x \)}{\( \frac{\sin x}{1-\sin x} \)}{\( \mathrm{tg}^2 x \)}{\( \mathrm{tg}\,2x \)}{}{}}
\LANGes{
\Question{
La expresión \( \frac{\sin 2x}{\cos^2x } \) equivale a:}
{\( 2\,\mathrm{tg}\,x \)}{\( \frac{\sin x}{1-\sin x} \)}{\( \mathrm{tg}^2 x \)}{\( \mathrm{tg}\,2x \)}{}{}}
\End

\Data\ID{1003076704}
\Updated{2020-07-15 03:07:12}
\Level{C}
\SubArea{Sine, cosine, tangent and cotangent}
\LANGen{
\Question{
Simplifying the expression \( 1 - (\cos x - \sin x)^2 \) we get:}
{\( \sin 2x \)}{\( \cos 2x \)}{\( 2 - \sin x \)}{\( 1 - \sin 2x \)}{}{}}
\LANGpl{
\Question{
Uprość wyrażenie \( 1 - (\cos x - \sin x)^2 \).}
{\( \sin 2x \)}{\( \cos 2x \)}{\( 2 - \sin x \)}{\( 1 - \sin 2x \)}{}{}}
\LANGcs{
\Question{
Zjednodušením výrazu \( 1 - (\cos x - \sin x)^2 \) dostaneme:}
{\( \sin 2x \)}{\( \cos 2x \)}{\( 2 - \sin x \)}{\( 1 - \sin 2x \)}{}{}}
\LANGsk{
\Question{
Po úprave výrazu   \( 1 - (\cos x - \sin x)^2 \) dostaneme:}
{\( \sin 2x \)}{\( \cos 2x \)}{\( 2 - \sin x \)}{\( 1 - \sin 2x \)}{}{}}
\LANGes{
\Question{
Al simplificar la expresión \( 1 - (\cos x - \sin x)^2 \) obtenemos:}
{\( \sin 2x \)}{\( \cos 2x \)}{\( 2 - \sin x \)}{\( 1 - \sin 2x \)}{}{}}
\End

\Data\ID{1003076705}
\Updated{2020-07-15 03:07:12}
\Level{C}
\SubArea{Sine, cosine, tangent and cotangent}
\LANGen{
\Question{
At which point \( x \in [12\pi;14\pi] \) has the function \( y = \sin x \) a maximum?}
{\( 12.5\pi \)}{\( 13\pi \)}{\( 12\pi \)}{\( 13.5\pi \)}{}{}}
\LANGpl{
\Question{
W którym punkcie \( x \in \langle12\pi;14\pi\rangle \) funkcja \( y = \sin x \) ma maksimum?}
{\( 12{,}5\pi \)}{\( 13\pi \)}{\( 12\pi \)}{\( 13{,}5\pi \)}{}{}}
\LANGcs{
\Question{
Pro které \( x \in \langle12\pi;14\pi\rangle \)  nabývá funkce \( y = \sin x \) maxima?}
{\( 12{,}5\pi \)}{\( 13\pi \)}{\( 12\pi \)}{\( 13{,}5\pi \)}{}{}}
\LANGsk{
\Question{
Pre ktoré \( x \in \langle12\pi;14\pi\rangle \)   nadobúda funkcia \( y = \sin x \)   maximum?}
{\( 12{,}5\pi \)}{\( 13\pi \)}{\( 12\pi \)}{\( 13{,}5\pi \)}{}{}}
\LANGes{
\Question{
En qué punto \( x \in [12\pi;14\pi] \) tiene la función \( y = \sin x \) su máximo?}
{\( 12.5\pi \)}{\( 13\pi \)}{\( 12\pi \)}{\( 13.5\pi \)}{}{}}
\End

\Data\ID{1003076708}
\Updated{2020-07-15 03:07:12}
\Level{C}
\SubArea{Sine, cosine, tangent and cotangent}
\LANGen{
\Question{
The measures of the interior angles of a  triangle are \( 30^{\circ} \), \( 45^{\circ} \) and \( 105^{\circ} \). The length of its longest side is \( 10\,\mathrm{cm} \). The length of its shortest side is:}
{\( 5.18\,\mathrm{cm} \)}{\( 7.33\,\mathrm{cm} \)}{\( 5.01\,\mathrm{cm} \)}{\( 7.07\,\mathrm{cm} \)}{}{}}
\LANGpl{
\Question{
Miary kątów wewnętrznych trójkąta są równe  \( 30^{\circ} \), \( 45^{\circ} \) i \( 105^{\circ} \). Długość najdłuższego boku wynosi \( 10\,\mathrm{cm} \). Długość najkrótszego boku jest równa:}
{\( 5{,}18\,\mathrm{cm} \)}{\( 7{,}33\,\mathrm{cm} \)}{\( 5{,}01\,\mathrm{cm} \)}{\( 7{,}07\,\mathrm{cm} \)}{}{}}
\LANGcs{
\Question{
Vnitřní úhly trojúhelníku mají velikost \( 30^{\circ} \), \( 45^{\circ} \) a \( 105^{\circ} \). Jeho nejdelší strana měří \( 10\,\mathrm{cm} \). Nejkratší strana trojúhelníku měří:}
{\( 5{,}18\,\mathrm{cm} \)}{\( 7{,}33\,\mathrm{cm} \)}{\( 5{,}01\,\mathrm{cm} \)}{\( 7{,}07\,\mathrm{cm} \)}{}{}}
\LANGsk{
\Question{
Vnútorné uhly trojuholníka majú veľkosti  \( 30^{\circ} \), \( 45^{\circ} \) a \( 105^{\circ} \). Jeho najdlhšia strana meria  \( 10\,\mathrm{cm} \). Najkratšia strana trojuholníka meria:}
{\( 5{,}18\,\mathrm{cm} \)}{\( 7{,}33\,\mathrm{cm} \)}{\( 5{,}01\,\mathrm{cm} \)}{\( 7{,}07\,\mathrm{cm} \)}{}{}}
\LANGes{
\Question{
Las medidas de los ángulos interiores de un triángulo son  \( 30^{\circ} \), \( 45^{\circ} \) y \( 105^{\circ} \).  La longitud de su lado más largo es\( 10\,\mathrm{cm} \). La longitud de su lado más corto es:}
{\( 5.18\,\mathrm{cm} \)}{\( 7.33\,\mathrm{cm} \)}{\( 5.01\,\mathrm{cm} \)}{\( 7.07\,\mathrm{cm} \)}{}{}}
\End

\Data\ID{1003076709}
\Updated{2020-07-15 03:07:12}
\Level{C}
\SubArea{Sine, cosine, tangent and cotangent}
\LANGen{
\Question{
The area of an obtuse triangle is \( 2\,\mathrm{dm}^2 \). The lengths of sides containing the obtuse angle are \( 2\,\mathrm{dm} \) and \( 4\,\mathrm{dm} \). The measure of this  angle is:}
{\( 150^{\circ} \)}{\( 165^{\circ} \)}{\( 155^{\circ} \)}{\( 158^{\circ} \)}{}{}}
\LANGpl{
\Question{
Pole powierzchni trójkąta rozwartokątnego wynosi \( 2\,\mathrm{dm}^2 \). Długości boków zawierających kąt rozwarty są równe \( 2\,\mathrm{dm} \) i \( 4\,\mathrm{dm} \). Wskaż miarę tego kąta:}
{\( 150^{\circ} \)}{\( 165^{\circ} \)}{\( 155^{\circ} \)}{\( 158^{\circ} \)}{}{}}
\LANGcs{
\Question{
Tupoúhlý trojúhelník má obsah \( 2\,\mathrm{dm}^2 \). Strany svírající tupý úhel jsou dlouhé \( 2\,\mathrm{dm} \) a \( 4\,\mathrm{dm} \). Velikost tupého úhlu v trojúhelníku je:}
{\( 150^{\circ} \)}{\( 165^{\circ} \)}{\( 155^{\circ} \)}{\( 158^{\circ} \)}{}{}}
\LANGsk{
\Question{
Tupouhlý trojuholník má obsah  \( 2\,\mathrm{dm}^2 \). Strany určujúce tupý uhol sú dlhé  \( 2\,\mathrm{dm} \) a \( 4\,\mathrm{dm} \). Veľkosť tupého uhla v trojuholníku je:}
{\( 150^{\circ} \)}{\( 165^{\circ} \)}{\( 155^{\circ} \)}{\( 158^{\circ} \)}{}{}}
\LANGes{
\Question{
El área de un triángulo obtuso es  \( 2\,\mathrm{dm}^2 \).  Las longitudes de los lados que contienen el ángulo obtuso son \( 2\,\mathrm{dm} \) y \( 4\,\mathrm{dm} \). La medida de este ángulo es:}
{\( 150^{\circ} \)}{\( 165^{\circ} \)}{\( 155^{\circ} \)}{\( 158^{\circ} \)}{}{}}
\End

\Data\ID{1003076710}
\Updated{2020-07-15 03:07:12}
\Level{C}
\SubArea{Sine, cosine, tangent and cotangent}
\LANGen{
\Question{
\( ABC \) is a triangle where the side \( b \) is \( 74\,\mathrm{cm} \) long and the angle \( \alpha = 60^{\circ} \). Find the length of its side \( c \)  if you know that the area of the triangle is \( 720.9\,\mathrm{cm}^2 \).}
{\( 22.5\,\mathrm{cm} \)}{\( 37.56\,\mathrm{cm} \)}{\( 38.97\,\mathrm{cm} \)}{\( 24.54\,\mathrm{cm} \)}{}{}}
\LANGpl{
\Question{
Dany jest trójkąt \( ABC \) o boku  \( b \) równym \( 74\,\mathrm{cm} \) i kącie \( \alpha = 60^{\circ} \). Oblicz długość boku \( c \)  wiedząc, że pole powierzchni trójkąta wynosi \( 720{,}9\,\mathrm{cm}^2 \).}
{\( 22{,}5\,\mathrm{cm} \)}{\( 37{,}56\,\mathrm{cm} \)}{\( 38{,}97\,\mathrm{cm} \)}{\( 24{,}54\,\mathrm{cm} \)}{}{}}
\LANGcs{
\Question{
Jaká je velikost strany c v trojúhelníku \( ABC \), jestliže jeho obsah je  \( 720{,}9\,\mathrm{cm}^2 \), délka strany \( b \) je \( 74\,\mathrm{cm} \) a úhel \( \alpha = 60^{\circ} \)?}
{\( 22{,}5\,\mathrm{cm} \)}{\( 37{,}56\,\mathrm{cm} \)}{\( 38{,}97\,\mathrm{cm} \)}{\( 24{,}54\,\mathrm{cm} \)}{}{}}
\LANGsk{
\Question{
Aká je veľkosť strany c v trojuholníku \( ABC \), ak jeho obsah je  \( 720{,}9\,\mathrm{cm}^2 \), dĺžka strany \( b \) je \( 74\,\mathrm{cm} \) a uhol \( \alpha = 60^{\circ} \)?}
{\( 22{,}5\,\mathrm{cm} \)}{\( 37{,}56\,\mathrm{cm} \)}{\( 38{,}97\,\mathrm{cm} \)}{\( 24{,}54\,\mathrm{cm} \)}{}{}}
\LANGes{
\Question{
\( ABC \) es un triángulo donde el lado \( b \)  es  \( 74\,\mathrm{cm} \) largo y el ángulo \( \alpha = 60^{\circ} \). Calcula la longitud de su lado \( c \)  si sabes que el área del triángulo es \( 720.9\,\mathrm{cm}^2 \).}
{\( 22.5\,\mathrm{cm} \)}{\( 37.56\,\mathrm{cm} \)}{\( 38.97\,\mathrm{cm} \)}{\( 24.54\,\mathrm{cm} \)}{}{}}
\End

\Data\ID{1003076711}
\Updated{2020-07-15 03:07:12}
\Level{C}
\SubArea{Sine, cosine, tangent and cotangent}
\LANGen{
\Question{
\( ABC \) is a triangle with side \( c = 27\,\mathrm{cm} \) and the angle \( \gamma = 78^{\circ} \). The radius of the circle  circumscribed is:}
{\( 13.8\,\mathrm{cm} \)}{\( 27.6\,\mathrm{cm} \)}{\( 18.3\,\mathrm{cm} \)}{\( 6.9\,\mathrm{cm} \)}{}{}}
\LANGpl{
\Question{
Dany jest trójką \( ABC \) o boku \( c = 27\,\mathrm{cm} \) i kącie  \( \gamma = 78^{\circ} \). Promień okręgu opisanego wynosi:}
{\( 13{,}8\,\mathrm{cm} \)}{\( 27{,}6\,\mathrm{cm} \)}{\( 18{,}3\,\mathrm{cm} \)}{\( 6{,}9\,\mathrm{cm} \)}{}{}}
\LANGcs{
\Question{
V trojúhelníku \( ABC \)  je dána strana \( c = 27\,\mathrm{cm} \) a úhel \( \gamma = 78^{\circ} \). Poloměr kružnice opsané je:}
{\( 13{,}8\,\mathrm{cm} \)}{\( 27{,}6\,\mathrm{cm} \)}{\( 18{,}3\,\mathrm{cm} \)}{\( 6{,}9\,\mathrm{cm} \)}{}{}}
\LANGsk{
\Question{
V trojuholníku \( ABC \)  je dané: strana  \( c = 27\,\mathrm{cm} \) a uhol \( \gamma = 78^{\circ} \). Polomer opísanej kružnice je:}
{\( 13{,}8\,\mathrm{cm} \)}{\( 27{,}6\,\mathrm{cm} \)}{\( 18{,}3\,\mathrm{cm} \)}{\( 6{,}9\,\mathrm{cm} \)}{}{}}
\LANGes{
\Question{
\( ABC \) es un triángulo con el lado \( c = 27\,\mathrm{cm} \) y el ángulo \( \gamma = 78^{\circ} \).  El radio de la circunferencia circunscrita es:}
{\( 13.8\,\mathrm{cm} \)}{\( 27.6\,\mathrm{cm} \)}{\( 18.3\,\mathrm{cm} \)}{\( 6.9\,\mathrm{cm} \)}{}{}}
\End

\Data\ID{1003076712}
\Updated{2020-07-15 03:07:12}
\Level{C}
\SubArea{Sine, cosine, tangent and cotangent}
\LANGen{
\Question{
Let \( \sin x = a \) and \( \cos x = b \). Then \( \cos(2x) \) is equal to:}
{\( b^2-a^2 \)}{\( 2b \)}{\( 2ab \)}{\( a^2-b^2 \)}{}{}}
\LANGpl{
\Question{
Jeśli \( \sin x = a \) i \( \cos x = b \), wtedy \( \cos(2x) \) jest równy:}
{\( b^2-a^2 \)}{\( 2b \)}{\( 2ab \)}{\( a^2-b^2 \)}{}{}}
\LANGcs{
\Question{
Jestliže \( \sin x = a \) a \( \cos x = b \), pak můžeme vyjádřit \( \cos(2x) \)  jako:}
{\( b^2-a^2 \)}{\( 2b \)}{\( 2ab \)}{\( a^2-b^2 \)}{}{}}
\LANGsk{
\Question{
Ak \( \sin x = a \) a \( \cos x = b \), tak \( \cos(2x) \) vieme vyjadriť ako:}
{\( b^2-a^2 \)}{\( 2b \)}{\( 2ab \)}{\( a^2-b^2 \)}{}{}}
\LANGes{
\Question{
Sea \( \sin x = a \) y \( \cos x = b \). Entonces \( \cos(2x) \) equivale a:}
{\( b^2-a^2 \)}{\( 2b \)}{\( 2ab \)}{\( a^2-b^2 \)}{}{}}
\End

\Data\ID{1103076608}
\Updated{2020-07-15 03:07:12}
\Level{C}
\SubArea{Sine, cosine, tangent and cotangent}
\IMG{1103076608_0.pdf}
\LANGen{
\Question{
The figure shows the graph of the function:}
{\( f(x) = -\sin(-2x) \)}{\( f(x) = -\sin 2x \)}{\( f(x) = |\sin 2x| \)}{\( f(x) = \sin|2x| \)}{}{}}
\LANGpl{
\Question{
Rysunek przedstawia wykres funkcji:}
{\( f(x) = -\sin(-2x) \)}{\( f(x) = -\sin 2x \)}{\( f(x) = |\sin 2x| \)}{\( f(x) = \sin|2x| \)}{}{}}
\LANGcs{
\Question{
Na obrázku je část grafu funkce:}
{\( f(x) = -\sin(-2x) \)}{\( f(x) = -\sin 2x \)}{\( f(x) = |\sin 2x| \)}{\( f(x) = \sin|2x| \)}{}{}}
\LANGsk{
\Question{
Na obrázku je časť grafu funkcie:}
{\( f(x) = -\sin(-2x) \)}{\( f(x) = -\sin 2x \)}{\( f(x) = |\sin 2x| \)}{\( f(x) = \sin|2x| \)}{}{}}
\LANGes{
\Question{
La figura muestra la gráfica de la función:}
{\( f(x) = -\sin(-2x) \)}{\( f(x) = -\sin 2x \)}{\( f(x) = |\sin 2x| \)}{\( f(x) = \sin|2x| \)}{}{}}
\End

\Data\ID{1103076610}
\Updated{2020-07-15 03:07:12}
\Level{C}
\SubArea{Sine, cosine, tangent and cotangent}
\IMG{1103076610.pdf}
\LANGen{
\Question{
The figure shows the graph of the function:}
{\( f(x)=\cos(-2x) \)}{\( f(x)=-\cos 2x \)}{\( f(x)=|\cos 2x| \)}{\( f(x)=-|\cos 2x|   \)}{}{}}
\LANGpl{
\Question{
Rysunek przedstawia wykres funkcji:}
{\( f(x)=\cos(-2x) \)}{\( f(x)=-\cos 2x \)}{\( f(x)=|\cos 2x| \)}{\( f(x)=-|\cos 2x|   \)}{}{}}
\LANGcs{
\Question{
Na obrázku je graf funkce:}
{\( f(x)=\cos(-2x) \)}{\( f(x)=-\cos 2x \)}{\( f(x)=|\cos 2x| \)}{\( f(x)=-|\cos 2x|   \)}{}{}}
\LANGsk{
\Question{
Na obrázku je časť grafu funkcie:}
{\( f(x)=\cos(-2x) \)}{\( f(x)=-\cos 2x \)}{\( f(x)=|\cos 2x| \)}{\( f(x)=-|\cos 2x|   \)}{}{}}
\LANGes{
\Question{
La figura muestra la gráfica de la función:}
{\( f(x)=\cos(-2x) \)}{\( f(x)=-\cos 2x \)}{\( f(x)=|\cos 2x| \)}{\( f(x)=-|\cos 2x|   \)}{}{}}
\End

\Data\ID{1103076611}
\Updated{2020-07-15 03:07:12}
\Level{C}
\SubArea{Sine, cosine, tangent and cotangent}
\IMG{1103076611.pdf}
\LANGen{
\Question{
The graph shown in the picture is graph of  the function:}
{\( f(x)=\sin|x| \)}{\( f(x)=|\sin x| \)}{\( f(x)=|\cos x| \)}{\( f(x)=\cos|x| \)}{}{}}
\LANGpl{
\Question{
Wskaż funkcję, której wykres jest przedstawiony na rysunku:}
{\( f(x)=\sin|x| \)}{\( f(x)=|\sin x| \)}{\( f(x)=|\cos x| \)}{\( f(x)=\cos|x| \)}{}{}}
\LANGcs{
\Question{
Na obrázku je část grafu funkce:}
{\( f(x)=\sin|x| \)}{\( f(x)=|\sin x| \)}{\( f(x)=|\cos x| \)}{\( f(x)=\cos|x| \)}{}{}}
\LANGsk{
\Question{
Na obrázku je časť grafu funkcie:}
{\( f(x)=\sin|x| \)}{\( f(x)=|\sin x| \)}{\( f(x)=|\cos x| \)}{\( f(x)=\cos|x| \)}{}{}}
\LANGes{
\Question{
La gráfica mostrada en la imagen es la gráfica de la función:}
{\( f(x)=\sin|x| \)}{\( f(x)=|\sin x| \)}{\( f(x)=|\cos x| \)}{\( f(x)=\cos|x| \)}{}{}}
\End

\Data\ID{1103082703}
\Updated{2020-07-15 03:07:12}
\Level{C}
\SubArea{Sine, cosine, tangent and cotangent}
\IMG{1103082703_0.pdf}
\LANGen{
\Question{
Function \( f \) is given completely by the next graph. Identify which of the following statements is true.}
{\( f(x)=-|\sin x|;\ x\in[-2\pi;2\pi] \)}{\( f(x)=|\cos x|;\ x\in[-2\pi;2\pi] \)}{\( f(x)=|-\sin x|;\ x\in[-2\pi;2\pi] \)}{\( f(x)=-0.5\cdot\sin x;\ x\in[-2\pi;2\pi] \)}{}{}}
\LANGpl{
\Question{
Wykres przedstawia funkcję \( f \). Wskaż, które stwierdzenie jest prawdziwe.}
{\( f(x)=-|\sin x|;\ x\in\langle-2\pi;2\pi\rangle \)}{\( f(x)=|\cos x|;\ x\in\langle-2\pi;2\pi\rangle \)}{\( f(x)=|-\sin x|;\ x\in\langle-2\pi;2\pi\rangle \)}{\( f(x)=-0{,}5\cdot\sin x;\ x\in\langle-2\pi;2\pi\rangle \)}{}{}}
\LANGcs{
\Question{
Funkce \( f \) je dána následujícím grafem. Určete, které z následujících tvrzení je pravdivé.}
{\( f(x)=-|\sin x|;\ x\in\langle-2\pi;2\pi\rangle \)}{\( f(x)=|\cos x|;\ x\in\langle-2\pi;2\pi\rangle \)}{\( f(x)=|-\sin x|;\ x\in\langle-2\pi;2\pi\rangle \)}{\( f(x)=-0{,}5\cdot\sin x;\ x\in\langle-2\pi;2\pi\rangle \)}{}{}}
\LANGsk{
\Question{
Funkcia \( f \) je daná následujúcim grafom. Určte, ktoré z následujúcich tvrdení je pravdivé.}
{\( f(x)=-|\sin x|;\ x\in\langle-2\pi;2\pi\rangle \)}{\( f(x)=|\cos x|;\ x\in\langle-2\pi;2\pi\rangle \)}{\( f(x)=|-\sin x|;\ x\in\langle-2\pi;2\pi\rangle \)}{\( f(x)=-0{,}5\cdot\sin x;\ x\in\langle-2\pi;2\pi\rangle \)}{}{}}
\LANGes{
\Question{
La función \( f \)  viene dada completamente por la siguiente gráfica. Identifica cuál de las siguientes proposiciones es verdadera.}
{\( f(x)=-|\sin x|;\ x\in[-2\pi;2\pi] \)}{\( f(x)=|\cos x|;\ x\in[-2\pi;2\pi] \)}{\( f(x)=|-\sin x|;\ x\in[-2\pi;2\pi] \)}{\( f(x)=-0.5\cdot\sin x;\ x\in[-2\pi;2\pi] \)}{}{}}
\End

\Data\ID{9000033810}
\Updated{2020-07-15 03:07:34}
\Level{C}
\SubArea{Sine, cosine, tangent and cotangent}
\LANGen{
\Question{
Find the parity (odd/even) of the function
\(l\colon y = |\mathop{\mathrm{cotg}}\nolimits x|\).}
{The function \(l\) 
is an even function.}{The function \(l\) 
is an odd function.}{The function \(l\) 
is neither an odd nor even function.}{}{}{}}
\LANGpl{
\Question{
Sprawdź parzystość (tzn.parzysta/nieparzysta) funkcji
\(l\colon y = |\mathop{\mathrm{cotg}}\nolimits x|\).}
{Funkcja \(l\) 
jest parzysta.}{Funkcja \(l\) 
jest nieparzysta.}{Funkcja \(l\) 
nie jest parzysta, ani nieparzysta.}{}{}{}}
\LANGcs{
\Question{
Rozhodněte o paritě (tzn. o sudosti/lichosti) funkce
\(l\colon y = |\mathop{\mathrm{cotg}}\nolimits x|\).}
{Funkce \(l\) 
je sudá.}{Funkce \(l\) 
je lichá.}{Funkce \(l\) 
není ani sudá, ani lichá.}{}{}{}}
\LANGsk{
\Question{
Rozhodnite o párnosti alebo nepárnosti funkcie
\(l\colon y = |\mathop{\mathrm{cotg}}\nolimits x|\).}
{Funkcia \(l\) 
je párna.}{Funkcia \(l\) 
je nepárna.}{Funkcia \(l\) 
nie je ani párna, ani nepárna.}{}{}{}}
\LANGes{
\Question{
Determina la paridad (impar / par) de la función
\(l\colon y = |\mathop{\mathrm{cotg}}\nolimits x|\).}
{La función  \(l\) 
es una función par.}{La función   \(l\) 
es una función impar.}{La función  \(l\) 
no es una función impar ni par.}{}{}{}}
\End
